\documentclass{amsart}
\title{Analysis}
\author{}

% Include packages
\usepackage{amsfonts}
\usepackage{amsmath}
\usepackage{amssymb}
\usepackage{amsthm}
% \usepackage{hyperref}
\usepackage{mathrsfs}

% Options
% \mathtoolsset{showonlyrefs}
% \allowdisplaybreaks
\numberwithin{equation}{section}

% Begin Definitions
\DeclareMathOperator{\mesh}{mesh}
\newcommand{\dif}{\mathrm{d}}
\newcommand{\Dif}{\mathrm{D}}
\newcommand{\lequiv}{:\Leftrightarrow}
\newcommand{\lxor}{\otimes}
\newcommand{\comp}{\backslash}
\newcommand{\union}{\cup}
% Custom Proof environment to allow cases
\expandafter\let\expandafter\oldproof\csname\string\proof\endcsname
\let\oldendproof\endproof
\renewenvironment{proof}[1][\proofname]{%
  \setcounter{case}{0}
  \setcounter{step}{0}
  \oldproof[#1]%
}{\oldendproof}
% Changes the header to \partmark{Part Title}
\newcommand{\partmark}[1]{\markboth{\MakeUppercase{#1}}{\MakeUppercase{#1}}}
% Theorem
\theoremstyle{plain}
\newtheorem*{theorem*}{Theorem}
\newtheorem{theorem}{Theorem}[section]
\newtheorem{axiom}[theorem]{Axiom}
\newtheorem{corollary}[theorem]{Corollary}
\newtheorem{lemma}[theorem]{Lemma}
\theoremstyle{definition}
\newtheorem{definition}[theorem]{Definition}
\newtheorem*{definition*}{Definition}
\theoremstyle{remark}
\newtheorem*{remark}{Remark}
\newtheorem*{example}{Example}
\newtheorem{case}{Case}
\newtheorem*{case*}{Case}
\newtheorem{step}{Step}

% Begin Document
\begin{document}
\maketitle
% \setcounter{tocdepth}{1}
% \tableofcontents
\partmark{Real Analysis}
\section{Convergent Sequences}
\begin{definition}[Metric]
  Let $M$ be a set and $R$ be an ordered ring. A \emph{metric} on $M$ is a function
  \begin{align*}
    d: M \times M \to R
  \end{align*}
  satisfying, for all $x,y,z\in M$:
  \begin{enumerate}
    \item \emph{Non-negativity}: $d(x,y) \geq 0$,
    \item \emph{Identity of indiscernibles}: $d(x,y) = 0 \iff x = y$,
    \item \emph{Symmetry}: $d(x,y) = d(y,x)$,
    \item \emph{Triangle inequality}: $d(x,z) \leq d(x,y) + d(y,z)$.
  \end{enumerate}
  A \emph{metric space} $\mathscr{M}$ is a pair $(M, d)$ where $M$ is a set and $d$ is a metric on $M$.
\end{definition}

\begin{definition}[Bounded Sequence]
  A sequence $(s_n)$ in a metric space $(M, d)$ is \emph{bounded} if there exists $p \in M$ and $r>0$ such that
  \begin{align*}
    d(s_n, p) \leq r \quad \text{for all } n \in \mathbb{N}.
  \end{align*}
\end{definition}

\begin{definition}[Convergent Sequence]
  A sequence $(s_n)$ in a metric space $(M, d)$ \emph{converges} to $l \in M$ if, for every $\varepsilon>0$, there exists an integer $N \in \mathbb{N}$ such that whenever $n > N$, we have:
  \begin{align*}
    d(s_n, l) < \varepsilon.
  \end{align*}
  The limit of a sequence is denoted by
  \begin{align*}
    l = \lim_n s_n.
  \end{align*}
\end{definition}

\begin{theorem}\label{Thm:ConvergentImpliesBounded}
  Every convergent sequence is bounded.
\end{theorem}
\begin{proof}
  Let $(s_n)$ be a sequence in a metric space $(M, d)$ that converges to $l\in M$. Since $(s_n)$ converges, there exists $\varepsilon > 0$ and some $N \in \mathbb{N}$ such that for all $n > N$, we have
  \begin{align*}
    d(s_n, l) < \varepsilon.
  \end{align*}
  The finitely many terms $\{s_n : 0 \leq n \leq N\}$ are each at some distance from $l$, so define
  \begin{align*}
    B' := \max_{0 \leq n \leq N} \{ d(s_n, l) \}.
  \end{align*}
  Now let
  \begin{align*}
    B := \max \{B', \varepsilon\}.
  \end{align*}
  Then for all $n\in\mathbb{N}$, we have $d(s_n, l) \leq B$, so the entire sequence lies within the ball of radius $B$ centered at $l$. Hence, $s_n$ is bounded.
\end{proof}

\begin{definition}[Commutative Metric Group]
  A triple $(M, +, d)$ is called a \emph{commutative metric group} if:
  \begin{enumerate}
    \item $(M, +)$ forms a commutative group, and
    \item $(M, d)$ forms a metric space.
  \end{enumerate}
  In these notes, whenever we work with such a structure using the symbol $+$, we tacitly assume the metric is \emph{translation-invariant} with respect to this operation:
  \begin{align*}
    d(x + z, y + z) = d(x, y)\text{ for all }x, y, z \in M.
  \end{align*}
\end{definition}

\begin{theorem}
  Let $(s_n)$ and $(t_n)$ be convergent sequences in a commutative metric group $(M, +, d)$. Then the sequence $(s_n + t_n)$ also converges, and
  \begin{align*}
    \lim_n (s_n + t_n) &= \lim_n s_n + \lim_n t_n.
  \end{align*}
\end{theorem}
\begin{proof}
  Let $(s_n)$ and $(t_n)$ converge to $S$ and $T$ respectively.
  Choose $\varepsilon'>0$ and define $\varepsilon:= 2\varepsilon'$.
  By the definition of convergence, there exists integers $N_s$ and $N_t$ such that:
  \begin{align*}
    n > N_s \implies d(s_n, S) < \varepsilon'\\
    n > N_t \implies d(t_n, T) < \varepsilon'.
  \end{align*}
  Let $N:= \max \{ N_s, N_t\}$, then for all $n>N$, we compute
  \begin{align*}
    d(s_n+t_n, S+T) &= d(s_n+t_n-S-T,0)&&\text{(by translation invariance)},\\
    &= d((s_n-S)+(t_n-T),0),\\
    &\leq d(s_n-S, 0) + d(t_n-T, 0)&&\text{(by the triangle identity)},\\
    &= d(s_n, S) + d(t_n, T)&&\text{(by translation invariance)}.
  \end{align*}
  By the definition of convergence
  \begin{align*}
    d(s_n+t_n, S+T) \leq d(s_n, S) + d(t_n, T) < \varepsilon ' + \varepsilon' = \varepsilon.
  \end{align*}
  Thus, again by the definition of convergence
  \begin{align*}
    \lim_n (s_n+t_n) = S+T = \lim_n s_n + \lim_n t_n.
  \end{align*}
\end{proof}

\begin{definition}[Commutative Metric Ring]
  A quadruple $(M, +, \cdot, d)$ is called a \emph{commutative metric ring} if:
  \begin{enumerate}
    \item $(M, +, \cdot)$ is a commutative ring,
    \item $(M, d)$ is a metric space.
  \end{enumerate}
  In these notes, for a commutative metric ring, we assume:
  \begin{itemize}
    \item The metric is \emph{translation-invariant}:
    \begin{align*}
      d(x + z, y + z) = d(x, y)\text{ for all } x, y, z \in M,
    \end{align*}
    \item Multiplication is \emph{uniformly compatible} with the metric: there exists a constant $C > 0$ such that
    \begin{align*}
      d(x \cdot y, x \cdot z) \leq C \cdot d(y, z)\text{ for all } x, y, z \in M.
    \end{align*}
  \end{itemize}
\end{definition}

\begin{theorem}
  Let $(s_n)$ and $(t_n)$ be convergent sequences in a commutative metric ring $(M, +, \cdot, d)$. Then the sequence $(s_n \cdot t_n)$ also converges, and
  \begin{align*}
    \lim_n (s_n \cdot t_n) &= \left(\lim_n s_n\right)\left( \lim_n t_n\right).
  \end{align*}
\end{theorem}
\begin{proof}
  Let $(s_n)$ and $(t_n)$ converge to $S$ and $T$ respectively.
  Choose $\varepsilon'>0$.
  By the definition of convergence, there exists integers $N_s$ and $N_t$ such that:
  \begin{align*}
    n > N_s \implies d(s_n, S) < \varepsilon',\\
    n > N_t \implies d(t_n, T) < \varepsilon'.
  \end{align*}
  Let $N := \max\{N_s, N_t\}$, then for all $n > N$, we compute:
  \begin{align*}
    d(s_n \cdot t_n, S \cdot T)
    &= d(s_n \cdot t_n - S \cdot T, 0) && \text{(translation invariance)}\\
    &= d(s_n \cdot t_n -s_n \cdot T + s_n \cdot T - S \cdot T, 0)\\
    &= d(s_n \cdot (t_n - T) + (s_n - S) \cdot T, 0)\\
    &= d(s_n \cdot (t_n - T), -(s_n - S) \cdot T) && \text{(translation invariance)}\\
    &\leq d(s_n \cdot (t_n - T), 0) + d((S-s_n) \cdot T, 0) && \text{(triangle inequality)}\\
    &= d(s_n \cdot t_n , s_n \cdot T) +d(S \cdot T, s_n\cdot T) && \text{(translation invariance)}
  \end{align*}
  By uniform compatibility, there is $C_1, C_2>0$:
  \begin{align*}
    d(s_n\cdot t_n, s_n\cdot T)&\leq C_1\cdot d(t_n, T) < C_1 \cdot \varepsilon'\\
    d(s_n\cdot T, S\cdot T)&\leq C_2\cdot d(s_n, S) < C_2 \cdot \varepsilon'
  \end{align*}
  Let $\varepsilon:=C_1\cdot \varepsilon' + C_2 \cdot \varepsilon'$. Then for all $n > N$:
  \begin{align*}
    d(s_n \cdot t_n, S \cdot T) < \varepsilon
  \end{align*}
  Therefore, by the definition of convergence,
  \begin{align*}
    \lim_n \left(s_n \cdot t_n\right)= S \cdot T = \left(\lim_n s_n\right)\left( \lim_n t_n\right).
  \end{align*}
\end{proof}

\begin{theorem}
  If a sequence converges, then its limit is unique.
\end{theorem}
\begin{proof}
  Suppose a sequence $(s_n)$ in a metric space $(M, d)$ converges to two distinct limits $L$ and $L'$.
  By the identity of indiscernibles, $d(L, L')>0$.
  Let $\varepsilon:=d(L,L')>0$.
  Then from the definition of convergence, there is a $N$ and $N'$ such that
  \begin{align*}
    n > N \implies d(s_n, L) < \frac{\varepsilon}{2}, \\
    n > N' \implies d(s_n, L') < \frac{\varepsilon}{2}.
  \end{align*}
  Let $M:= \max\{N, N'\}$. Then for all $n>M$, we have
  \begin{align*}
    d(s_n, L) + d(s_n, L') < \varepsilon.
  \end{align*}
  By the triangle inequality,
  \begin{align*}
     d(L, L') \leq d(L, s_n) + d(s_n, L') < \varepsilon = d(L, L'),
  \end{align*}
  which is a contradiction.
  Therefore, no such distinct $L$ and $L'$ can exist. The limit is unique.
\end{proof}


\begin{theorem}[Convergence by Proximity]
  Let $(x_n)$ and $(y_n)$ be sequences in a metric space $(M, d)$.
  If $(x_n)$ converges to $L \in M$ and $d(x_n, y_n)$ converges to zero, then $(y_n)$ also converges to $L$.
\end{theorem}
\begin{proof}
  Assume $(x_n)$ converges to $L \in M$, and $d(x_n, y_n)$ converges to $0$.
  Fix $\varepsilon > 0$.
  By convergence, there exists integers $N_1, N_2\in\mathbb{N}$,
  \begin{align*}
    n > N_1 \implies d(x_n, L) < \frac{\varepsilon}{2},\text{ and}\\
    n > N_2 \implies d(x_n, y_n) < \frac{\varepsilon}{2}.
  \end{align*}
  Let $N:= \max\{N_1, N_2\}$.
  For all $n > N$, by the triangle inequality,
  \begin{align*}
    d(y_n, L) \leq d(y_n, x_n) + d(x_n, L) < \frac{\varepsilon}{2} + \frac{\varepsilon}{2} = \varepsilon.
  \end{align*}
  Hence, by the definition of convergence, $(y_n)$ converges to $L$.
\end{proof}

\section{Closed Sets}

\begin{definition}[Closed Set]
  A subset $S\subset M$ is \emph{closed} in $M$ if every sequence $\langle s_n \rangle$ in $S$ that converges in $M$ also converges in $S$.
  \begin{align}
    \lim_n s_n \in M \implies \lim_n s_n \in S
  \end{align}
\end{definition}

%\begin{definition}[Open Set]
%A \emph{open set} in $M$ is a subset $S\subset M$ such that
%\begin{align}
%\forall x\in S, y\in M: \exists \varepsilon >0: d(x,y) < \varepsilon \implies y \in S
%\end{align}
%\end{definition}
%
%\begin{theorem}
%If $S$ is closed in $M$, then $M-S$ is open.
%\end{theorem}
%\begin{proof}
%Let $S$ be a closed set, and $L\in M-S$. Suppose that $M-S$ is not open. Then
%\begin{align}
%\forall \varepsilon' >0: d(p,L) < \varepsilon' \land p\notin M-S\nonumber\\
%\forall \varepsilon' >0: d(p,L) < \varepsilon' \land p\in S\nonumber\\
%\end{align}
%Define a sequence $s:\mathbb{N}\to S$ such that
%\begin{align}
%d(s_n , L) < \frac{1}{n}
%\end{align}
%Let $\varepsilon>0$, by Archimedean property, there is an $N$ such that
%\begin{align}
%\frac{1}{N} < \varepsilon
%\end{align}
%For all $n>N\implies \frac{1}{n} < \frac{1}{N}$ so for all $\varepsilon >0$ there is an $N$, such that if $n>N$
%\begin{align}
%d(s_n, L) < \varepsilon
%\end{align}
%Or
%\begin{align}
%\lim_n s_n = L
%\end{align}
%Since $S$ is closed, $L\in S$ which is a contradiction, therefor $M-S$ is open.
%\end{proof}
%
%
%\begin{theorem}
%If $S$ is open in $M$, then $M-S$ is closed.
%\end{theorem}
%\begin{proof}
%Let $S\subset M$ be a open set, and $s: \mathbb{N}\to M- S$ be a convergent sequence.
%\begin{align}
%L = \lim_n s_n
%\end{align}
%Suppose $L\in S$. Since $S$ is open, 
%\begin{align}
%\forall x\in M: \exists \varepsilon > 0 : d(x,L) < \varepsilon \implies x\in S
%\end{align}
%And since $s$ is convergent, there is an $N$, such that $n>N$
%\begin{align}
%d(s_n , L ) < \varepsilon
%\end{align}
%Then for $n>N: s_n \in S$, which contradicts the definition of $s$. Therefor $L\notin S$, $L\in M-S$, and $M-S$ is closed.
%\end{proof}

\begin{theorem}
  The real interval $[a,b]$ is closed.
  \begin{align}
    \forall s: \mathbb{N} \to [a,b]: \lim_n s_n = L \implies L \in [a,b]
  \end{align}
\end{theorem}
\begin{proof}
  Suppose not, then
  \begin{align}
    \exists s: \mathbb{N} \to [a,b]: \lim_n = L \land (L<a \lor b< L)
  \end{align}
  From the definition of a limit
  \begin{align}
    \forall \varepsilon > 0 : \exists N\in \mathbb{N}: n>N \implies L - \varepsilon < s_n < L + \varepsilon
  \end{align}
  \begin{case}[$b<L$]
    Then $L-b>0$. Let $\varepsilon = L -b$, then there is a $N_1$, such that $n>N_1$
    \begin{align}
      L - L + b = b < s_n
    \end{align}
    but the definition of the sequence restricts the domain to $s_n< b$, so we obtain a contradiction.
  \end{case}
  \begin{case}[$a<L$]
    Then $a-L>0$, let $\varepsilon = a-L$, which implies there is a $N_2$, such that $n>N_2$
    \begin{align}
      s_n < L + a - L = a
    \end{align}
    which contradicts the restriction on the domain $a<s_n$
  \end{case}
\end{proof}

\begin{theorem}[Inheritence Principle]
  If $S$ is closed in $M$ and $N$ is a subspace of $M$, then $N\cap S$ is closed in $N$.
\end{theorem}
\begin{proof}
  Let $N$ be a subspace of $M$ and $S$ some closed subset of $M$. Define a sequence $s:\mathbb{N} \cap S$ such that
  \begin{align}
    \exists L \in N: \forall \varepsilon >0: \exists N \in \mathbb{N}: \forall n \in \mathbb{N}: |s_{N+n}- L | < \varepsilon
  \end{align}
  Since $S$ is closed and $N\cap M$, $L\in S\cap N$. Hence $S\cap N $ is closed.
\end{proof}

\begin{corollary}
  If $N$ is closed in $M$, then $S$ is closed in $N$ if and only if $S$ is closed in $M$.
\end{corollary}


\section{Monotonic Sequences}


\begin{definition}[Monotonic Sequence]
  An \emph{increasing sequence} is a sequence where
  \begin{align}
    m > n \implies s_m > s_n
  \end{align}
  A \emph{strictly increasing sequence} is a sequence where
  \begin{align}
    m > n \implies s_m \geq s_n
  \end{align}
  Similarly, a \emph{decreasing sequence} is a sequence where
  \begin{align}
    m > n \implies s_m < s_n
  \end{align}
  And a \emph{strictly decreasing sequence} is a sequence where
  \begin{align}
    m > n \implies s_m \leq s_n
  \end{align}
  A \emph{monotonic sequence} is a sequence that is either increasing or decreasing.
\end{definition}


\begin{lemma}
  If $n:\mathbb{Z}^+\to \mathbb{Z}^+$ is a strictly increasing sequence ($i>j\implies n_i> n_j$), then for any positive integer $N$ 
  \begin{align}
    \exists i \in \mathbb{Z}^+: n_i > N
  \end{align}
\end{lemma}
\begin{proof}
  Hypothesis for induction:
  \begin{align}
    \forall i \in \mathbb{Z}^+: n_i \geq i
  \end{align}
  \begin{case}[Base]
    If $n_1\in\mathbb{Z}^+$, then $n_1 \geq 1$.
  \end{case}
  \begin{case}[Induction]
    Assume without loss of generality,
    \begin{align}
      n_i \geq i
    \end{align}
    Since $n$ is strictly increasing
    \begin{align}
      n_{i+1} > n_i
    \end{align}
    By trichotomy
    \begin{align}
      n_{i+1} > i
    \end{align}
    Then
    \begin{align}
      n_{i+1} \geq i+1
    \end{align}
  \end{case}
  Thus by induction, 
  \begin{align}
    \forall i \in \mathbb{Z}^+: n_i \geq i
  \end{align}
\end{proof}



\section{Subsequential Limits}

\begin{definition}[Subsequence]
  A \emph{subsequence} of a sequence $s$ is a sequence that is a composition $s \circ n$ where $n:\mathbb{Z}^+\to \mathbb{Z}^+$ is strictly increasing.
\end{definition}

\begin{definition}[Tail]
  The $N$-tail of a sequence $s:n\mapsto s_n$ is the subsequence $s_{N+n}$, where $N\in\mathbb{Z}^{\geq 0}$
\end{definition}


\begin{theorem}
  If a sequence converges to a limit, then all of its subsequential limits are equal to that limit.
\end{theorem}
\begin{proof}
  Consider a sequence $s$, that converges to $L$, and $s\circ n$ a subsequence of $s$.
  Fix $\varepsilon>0$. Since $s$ converges to $L$, there exists an $N$ such that
  \begin{align}
    n > N \implies |s_n - L | < \varepsilon
  \end{align}
  Let $I:=\inf\{i\in \mathbb{Z}^+: n_i >N \}$ (The set is not empty, by the lemma above $n_N\geq N$).
  Since $n_I\geq I>N$
  \begin{align}
    n > n_I \implies |s_n - L | < \varepsilon
  \end{align}
  From the definition of strictly increasing
  \begin{align}
    i> I \implies n_i > n_I
  \end{align}
  So
  \begin{align}
    i > I \implies n_i> n_I \implies |s_{n_i} - L | < \varepsilon
  \end{align}
  Or
  \begin{align}
    \lim_i s_{n_i} = L
  \end{align}
\end{proof}

\begin{theorem}
  If a sequence is bounded, then all subsequences have the same bounds.
\end{theorem}
\begin{proof}
  Let $s$ be a sequence bounded by $(b, B)$
  \begin{align}
    \forall n: b\leq s_n \leq B
  \end{align}
  Let $s\circ n$ be some subsequence, then
  \begin{align}
    \forall i: b\leq s_{n_i} \leq B
  \end{align}
\end{proof}

\begin{theorem}
  Every subsequence of a monotonic sequence is the same type of monotonic.
\end{theorem}
\begin{proof}
  Let $s$ be a sequence and $\sim$ be some order relation on the monotonic sequence such that
  \begin{align}
    m> n \implies s_m \sim s_n
  \end{align}
  Let $s\circ n$ be some subsequence, then
  \begin{align}
    i> j \implies n_i > n_j \implies s_{n_i} \sim s_{n_j}
  \end{align}
\end{proof}

\begin{lemma}\label{Lemma:PointPeak}
  If a tail of a sequence has a maximum, the parent sequence has a maximum.
  %If a subsequence has a maximum, the parent sequence has a maximum?
\end{lemma}
\begin{proof}
  Let $a:= \max \{s_n:n>N\}$, which exists by assumption.
  Let $b:= \max \{s_n:n\leq N\}$, which exists because the set is finite.
  If we let $c:= \max\{a,b\}$, then $c=\max\{s_n: n\in\mathbb{Z}^+\}$.
\end{proof}

\begin{theorem}[Point Peak Theorem]
  Every real sequence has a monotonic subsequence.
\end{theorem}
\begin{proof}
  Let $\langle s_n\rangle$ be any real sequence. Either every tail of $\langle s_n\rangle$ has a maximum
  \begin{align}
    \forall N\in\mathbb{Z}^{\geq 0},\exists m, \forall n: s_{N+m} \geq s_{N+n}
  \end{align}
  or there is a tail without a maximum.
  \begin{align}
    \exists N\in\mathbb{Z}^{\geq 0}, \forall m, \exists n: s_{N+m} < s_{N+n}
  \end{align}
  First suppose every tail of $\langle s\rangle$ has a maximum
  \begin{align}
    \forall N\in\mathbb{Z}^{\geq 0},\exists m, \forall n: s_{N+m} \geq s_{N+n}\nonumber
  \end{align}
  Define a sequence $k:\mathbb{Z}^+ \to \mathbb{Z}^+$
  \begin{align}
    k_1 &:= \min \{ m| \forall n: s_m \geq s_n\}\\
    k_{i+1} &:= \min \{ m | \forall n: s_{k_i+m} \geq s_{k_i+n} \} + k_i
  \end{align}
  The sequence $s\circ k$ is a subsequence because $k$ strictly increasing
  \begin{align}
    m \in \mathbb{Z}^+ \implies k_{i+1} > k_i
  \end{align}
  From the definition of $k_{i+1}$
  \begin{align}
    \forall i,n: s_{k_{i+1}} \geq s_{k_i+n}
  \end{align}
  Since $k$ is strictly increasing
  \begin{align}
    k_{i+2} > k_{i+1} > k_i
  \end{align}
  Let $n$ be the positive integer such that
  \begin{align}
    n:= k_{i+2} - k_i \in \mathbb{Z}^+
  \end{align}
  Applying this $n$ to the definition
  \begin{align}
    s_{k_{i+1}} \geq s_{k_i+k_{i+2}-k_i} = s_{k_{i+2}}
  \end{align}
  which shows that $s\circ k$ is decreasing.
  \begin{align}
    s_{k_i} \geq s_{k_{i+1}}
  \end{align}
  Now suppose there is a tail that does not have a maximum
  \begin{align}
    \exists N\in\mathbb{Z}^{\geq 0}, \forall m, \exists n: s_{N+m} < s_{N+n}\nonumber
  \end{align}
  Every tail of the tail without a maximum has no maximum, by lemma \ref{Lemma:PointPeak}.
  \begin{align}
    \forall M\in\mathbb{Z}^{\geq 0}, m, \exists n: s_{N+M+m} < s_{N+M+n}
  \end{align}
  This implies there is always a larger number further down the sequence.
  \begin{align}
    \forall m, \exists n: s_{N+m} < s_{N+m+n}
  \end{align}
  Define a sequence $k:\mathbb{Z}^+ \to \mathbb{Z}^+$
  \begin{align}
    k_1&:= N+1\\
    k_{i+1} &:= \min\{ n : s_{k_i+n} > s_{k_i} \} +k_i
  \end{align}
  The sequence $s\circ k$ is a subsequence because $k$ is strictly increasing
  \begin{align}
    n \in \mathbb{Z}^+ \implies k_{i+1} > k_i
  \end{align}
  From the definition of $k_{i+1}$, $s\circ k$ is strictly increasing
  \begin{align}
    s_{k_{i+1}} > s_{k_i}
  \end{align}
  Thus all real sequences have monotonic subsequences.
\end{proof}

\begin{theorem}
  For $|a|<1$
  \begin{align}
    \lim_n a^n = 0
  \end{align}
\end{theorem}
\begin{proof}
  Let $a$ be a number such that $|a| < 1$. Since $|a|>0$, multiplying the inequalities by $|a|$ gives
  \begin{align}
    0\geq|a|^2 &<|a|\\
    0\geq\left|a^2\right| &<|a|
  \end{align}
  Assume without loss of generality
  \begin{align}
    0\geq\left|a^{n+1}\right|< |a^n|
  \end{align}
  Then multiplying by $|a|$ gives
  \begin{align}
    0\geq\left|a^{n+1}\right||a|<\left|a^{n}\right||a|\\
    0\geq\left|a^{n+2}\right| < \left|a^{n+1}\right|
  \end{align}
  Then, by induction, for all $n\in\mathbb{Z}^+$
  \begin{align}
    0\geq\left|a^{n+1}\right| < \left|a^n\right|
  \end{align}
  which implies that the sequence $|a^n|$ is strictly decreasing, and bounded below by zero.
  By the monotonic convergence theorem, $|a^n|$ converges. Assume it converges to limit $L$.
  \begin{align}
    L:= \lim_n \left|a^n\right|
  \end{align}
  The sequence $\left|a^{2n}\right|$ is a subsequence of $\left|a^n\right|$, so it converges to the same limit
  \begin{align}
    L &= \lim_n \left|a^{2n}\right|\\
    &= \lim_n \left|a^n\right|^2\\
    &= \left(\lim_n \left|a^n\right|\right)^2\\
    &= L^2
  \end{align}
  Thus, either $L=0$ or $L=1$.
  Since $\left|a^1\right|=|a|<1$ and $a^n$ is strictly decreasing, $L \neq 1$. Therefor
  \begin{align}
    \lim_n |a^n| = 0
  \end{align}
  By lemma \ref{Lemma:ExponentialConvergentSequence}
  \begin{align}
    \lim_n a^n = 0
  \end{align}
\end{proof}

\section{Compact Sets}

\begin{definition}[Compact Set]
  A set is \emph{compact} if every sequence $s:\mathbb{Z}^+ \to S$ has a subsequential limit in $S$.
\end{definition}

\begin{theorem}
  A compact subset of $M$ is a closed subset of $M$
\end{theorem}
\begin{proof}
  Let $S$ be a compact subset of $M$ and $s: \mathbb{N} \to S$ be a sequence that converges to some $L\in M$. Since $s$ converges, every subsequential limit is equal to $L$. Since $S$ is compact $L\in S$. Thus $S$ is closed.
\end{proof}

\begin{theorem}
  A compact subset of $M$ is bounded in $M$
\end{theorem}
\begin{proof}
  Let $p$ be a point in $M$ and $S\subset M$. Suppose $S$ is not bounded, then there is some sequence $s:\mathbb{N}\to S$ defined by 
  \begin{align}
    d(s_n, p) \geq n
  \end{align}
  Since $S$ is compact, there is some subsequence of $s$ that converges. For every $\varepsilon>0$, there is an $N$
  \begin{align}
    i> N \implies d(s_{n_i} , L) < \varepsilon
  \end{align}
  Choose some $\varepsilon >0$. By the Archimedean property there is an $N'$ such that
  \begin{align}
    \varepsilon < N '
  \end{align}
  Let $M:= \max\{N,N'\}$. Since $M> N$
  \begin{align}
    d(s_{n_M} , p) < \varepsilon < N' < M
  \end{align}
  And since $s\circ n$ is a subsequence $M \leq n_M$, which along with the definition of the sequence gives
  \begin{align}
    M \leq n_M \leq d(s_{n_M} , p) < M
  \end{align}
  which is a contradiction. So $S$ is bounded.
\end{proof}


\section{Complete Spaces}

\begin{definition}[Pinched Sequence]
  A sequence $(s_n)$ in a metric space $(M, d)$ is \emph{pinched} if, for every $\varepsilon > 0$, there exists an integer $N \in \mathbb{N}$ such that whenever $n, m>N$, we have
  \begin{align*}
    d(s_n, s_m) < \varepsilon.
  \end{align*}
\end{definition}

\begin{definition}[Complete Space]
  A metric space $(M, d)$ is \emph{complete} if every pinched sequence in $M$ converges to a point in $M$.
  That is, for every sequence $(s_n)$ in $M$, if
  \begin{align*}
    \forall \varepsilon > 0\ \exists N \in \mathbb{N}\ \forall n, m> N: d(s_n, s_m) < \varepsilon,
  \end{align*}
  then there exists a point $L \in M$ such that
  \begin{align*}
    \forall \varepsilon > 0\ \exists N' \in \mathbb{N}\ \forall n > N': d(s_n, L) < \varepsilon.
  \end{align*}
\end{definition}

\begin{theorem}
  The metric space $(\mathbb{Q},d)$, where $d(x,y) = |x-y|$, is not complete.
\end{theorem}
\begin{proof}
  Proof by counter-example. Let $(F_n):\mathbb{N}\to\mathbb{N}$ be a sequence defined recursively:
  \begin{align*}
    F_{n+1} = F_n + F_{n-1}
  \end{align*}
  where $F_0=1$ and $F_1=1$. $(F_n)$ can be shown to be shown to be a sequence of increasing positive integers. Let $(x_n):\mathbb{N} \to \mathbb{Q}$ be a rational sequence defined by
  \begin{align*}
    x_n := \frac{F_{n+1}}{F_n}.
  \end{align*}
  Transforming $x_n$ into a recursive sequence:
  \begin{align*}
    x_{n+1} = \frac{F_{n+2}}{F_{n+1}} = \frac{F_{n+1}+F_n}{F_{n+1}} = 1 + \frac{F_n}{F_{n+1}} = 1+\frac{1}{x_n}.
  \end{align*}

  \step A well-known identity says that $F_{n+2}F_n-F_{n+1}^2=(-1)^n$. We'll show it by showing for the base case $(n=0)$:
  \begin{align*}
    (2)(1)-(1)^2=1=(-1)^0
  \end{align*}
  And then the inductive case $n=k \implies n=k+1$:
  \begin{align*}
    F_{k+3}F_{k+1}-F_{k+2}^2 &= (F_{k+2}+F_{k+1})F_{k+1}-F_{k+2}^2 \\
    &=F_{k+2}(F_{k+1}-F_{k+2})+F_{k+1}^2
  \end{align*}
  Using a re-arrangement of $F_{k+2}=F_{k+1}+F_k$:
  \begin{align*}
    F_{k+3}F_{k+1}-F_{k+2}^2 &= F_{k+2}(-F_k)+F_{k+1}^2\\
    &= (-1)(F_{k+2}F_k-F_{k+1}^2) = (-1)^{k+1}
  \end{align*}
  So by induction $F_{n+2}F_n-F_{n+1}^2=(-1)^n$ is true for all $n \in \mathbb{N}$.

  \step For all $n \geq 2$, we have $F_n\geq n$.
  We show this with a proof by induction.
  The base cases are $F_2=2\geq2$ and $F_3=3\geq3$.
  Now to show that $F_k \geq k$ and $F_{k-1}\geq k-1$ implies $F_{k+1}\geq k+1$:
  \begin{align*}
    F_{k+1} = F_k+F_{k-1} \geq k + (k-1) = 2k-1 \geq k+1.
  \end{align*}
  since $k\geq2$.
  \step Show that $|x_{n+1}-x_n|=\frac{1}{F_{n+1}F_n}$.
  Consider two successive terms $x_{n+1}$ and $x_n$:
  \begin{align*}
    |x_{n+1} - x_n| = \left|\frac{F_{n+2}}{F_{n+1}}-\frac{F_{n+1}}{F_n}\right| = \left|\frac{F_{n+2}F_n-F_{n+1}^2}{F_{n+1}F_n}\right|=\left| \frac{(-1)^n}{F_{n+1}F_n}\right|
  \end{align*}
  Since $F_n>0$ for all $n$, this can be further reduced to
  \begin{align*}
    |x_{n+1} - x_n| = \frac{1}{F_{n+1}F_n}
  \end{align*}
  \step We show that $(x_n)$ is pinched. Choose $\varepsilon>0$ and define
  \begin{align*}
    N:=\max\left\{2, \left\lceil\frac{1}{\varepsilon}\right\rceil\right\}.
  \end{align*}
  For $m>n>N$, the triangle inequality gives
  \begin{align*}
    |x_m-x_n| \leq \sum^{m-1}_{k=n} |x_{k+1}-x_k| = \sum^{m-1}_{k=n} \frac{1}{F_kF_{k+1}}
  \end{align*}
  Since $m>n\geq2$: $F_k \geq k$:
  \begin{align*}
    |x_m-x_n| \leq \sum^{m-1}_{k=n}\frac{1}{k(k+1)} &= \sum^{m-1}_{k=n} \left(\frac{1}{k}-\frac{1}{k+1}\right)\\
    &= \sum^{m-1}_{k=n}\frac{1}{k}-\sum^m_{k=n+1}\frac{1}{k} \\
    &= \frac{1}{n}+\sum^{m-1}_{k=n+1}\frac{1}{k}-\frac{1}{m}-\sum^{m-1}_{k=n+1}\frac{1}{k} \\
    &= \frac{1}{n} - \frac{1}{m} < \frac{1}{n} < \frac{1}{N}<\varepsilon.
  \end{align*}
  which implies that $(x_n)$ is pinched.
  \step Show by proof by contradiction that $(x_n)$ has no rational limit. Suppose a limit exists. By the recursive equation it must converge to an $L$ such that
  \begin{align*}
    L = 1 + \frac{1}{L}.
  \end{align*}
  Rearranging into a standard quadratic form:
  \begin{align*}
    L^2-L-1 = 0
  \end{align*}
  For a quadratic equation with integer coefficients, the roots are rational if and only if the discriminant is a perfect square.
  The discriminant for this equation is:
  \begin{align*}
    \Delta=(-1)^2-4(1)(-1)=5,
  \end{align*}
  which is not a perfect square, therefore $L$ is not rational.
\end{proof}

\begin{theorem}
  The sum of two fundamental sequences is fundamental.
\end{theorem}
\begin{proof}
  By the triangle inequality
  \begin{align}
    |s_m - s_n + t_m - t_n | \leq |s_m - s_n| + |t_m - t_n|
  \end{align}
  By the definition of fundamental sequences, $\forall \varepsilon_s, \varepsilon_t\in\mathbb{Q}^+$, there is a $N_s, N_t\in\mathbb{Z}^+$ and $N:=\max\{N_s, N_t\}$, such that if $m,n > N$
  \begin{align}
    |s_m - s_n| &< \varepsilon_s\nonumber\\
    |t_m - t_n| &< \varepsilon_t
  \end{align}
  By $a<b\land c<d\implies a+c<b+d$
  \begin{align}
    |s_m - s_n| + |t_m - t_n| < \varepsilon_s + \varepsilon_t
  \end{align}
  Let $\varepsilon:= \varepsilon_s+\varepsilon_t$
  \begin{align}
    |s_m - s_n +t_m - t_n| &< \varepsilon\\
    \left|(s_m + t_m) - (s_n + t_n) \right| &< \varepsilon
  \end{align}
\end{proof}

\begin{definition}[Equivalence]
  Two sequences are equivalent mod $\mathcal{Z}$ if and only if their difference is in $\mathcal{Z}$
  \begin{align}
    s\equiv_\mathcal{Z} t \iff s-t\in\mathcal{Z}
  \end{align}
\end{definition}

\begin{theorem}[Equivalence Reflexivity]
  \begin{align}
    a\equiv_\mathcal{Z} a 
  \end{align}
\end{theorem}
\begin{proof}
  \begin{align}
    a-a&: n \mapsto a_n - a_n = 0\\
    a-a&: n \mapsto 0\\
    a-a &= 0^*\in\mathbb{Z}
  \end{align}
  \begin{align}
    a\equiv_\mathcal{Z} a
  \end{align}
\end{proof}

\begin{theorem}[Equivalence Symmetry]
  \begin{align}
    a\equiv_\mathcal{Z} b \iff b\equiv_\mathcal{Z} a
  \end{align}
\end{theorem}
\begin{proof}
  \begin{align}
    a\equiv_\mathcal{Z} b
  \end{align}
  \begin{align}
    a-b&\in\mathcal{Z}\\
    -(a-b) &\in\mathcal{Z}\\
    b-a &\in\mathcal{Z}
  \end{align}
  \begin{align}
    b\equiv_\mathcal{Z} a
  \end{align}
\end{proof}

\begin{theorem}[Equivalence Transitivity]
  \begin{align}
    a\equiv_\mathcal{Z} b \land b \equiv_\mathcal{Z} c \implies a\equiv_\mathcal{Z} c
  \end{align}
\end{theorem}
\begin{proof}
  \begin{align}
    a\equiv_\mathcal{Z} b &\land b \equiv_\mathcal{Z} c\\
    a-b\in\mathcal{Z} &\land b-c \in\mathcal{Z}
  \end{align}
  Since addition is closed on $\mathbb{Z}$
  \begin{align}
    (a-b) + (b-c) &\in\mathcal{Z}\\
    a-c &\in\mathcal{Z}
  \end{align}
  \begin{align}
    a\equiv_\mathcal{Z} c
  \end{align}
\end{proof}

\begin{theorem}
  Addition is compatible with equivalence
  \begin{align}
    s\equiv_\mathcal{Z} s' \land t\equiv_\mathcal{Z} t' \implies s+t \equiv_\mathcal{Z} s'+t'
  \end{align}
\end{theorem}
\begin{proof}
  Since $s\equiv_\mathcal{Z} s'$ and $t\equiv_\mathcal{Z} t'$
  \begin{align}
    s-s'\in\mathcal{Z}\\
    t-t'\in\mathcal{Z}
  \end{align}
  Since $\mathcal{Z}$ is closed on addition
  \begin{align}
    (s-s')+(t-t')\in\mathcal{Z}\\
    (s+t)-(s'+t') \in\mathcal{Z}
  \end{align}
  which implies that
  \begin{align}
    s+t \equiv_\mathcal{Z} s'+t'
  \end{align}
\end{proof}

\begin{theorem}
  Multiplication is compatible with equivalence for bounded sequences
  \begin{align}
    \forall s,s',t,t'\in\mathcal{B}:s\equiv_\mathcal{Z} s' \land t\equiv_\mathcal{Z} t' \implies st \equiv_\mathcal{Z} s't'
  \end{align}
\end{theorem}
\begin{proof}
  Assume $s\equiv_\mathcal{Z} s'$ and $t\equiv_\mathcal{Z} t'$, then
  \begin{align}
    s-s'&\in\mathcal{Z}\\
    t-t'&\in\mathcal{Z}
  \end{align}
  Since $s\in\mathcal{B}$ and $t'\in\mathcal{B}$
  \begin{align}
    s(t-t')\in\mathcal{Z}\\
    st-st'\in\mathcal{Z}
  \end{align}
  and
  \begin{align}
    (s-s')t'\in\mathcal{Z}\\
    st'-s't'\in\mathcal{Z}
  \end{align}
  Since $\mathcal{Z}$ is closed on addition
  \begin{align}
    st-st'+st'-s't'\in\mathcal{Z}\\
    st-s't'\in\mathcal{Z}
  \end{align}
  By the definition of equivalence
  \begin{align}
    st \equiv_\mathcal{Z} s't'
  \end{align}
\end{proof}

\begin{theorem}
  If $x\equiv_\mathcal{z} y$ and $\lim_n x_n = L$, then $\lim_n y_n = L$
\end{theorem}
\begin{proof}
  Fix $\varepsilon\in\mathbb{Q}^+$
  \begin{align}
    \exists N_1\in\mathbb{Z}^+: \forall n> N_1: |x_n - y_n|< \varepsilon\\
    \exists N_2\in\mathbb{Z}^+: \forall n> N_2: |x_n - L| < \varepsilon
  \end{align}
  Let $N:= \max\{N_1, N_2\}$, if $n>N$
  \begin{align}
    |x_n-y_n| < \varepsilon\\
    |x_n - L| < \varepsilon
  \end{align}
  By the triangle identity
  \begin{align}
    |y_n-L | = |y_n-x_n+x_n-L| \leq |x_n-y_n| + |x_n- L| < \varepsilon +  \varepsilon = 2\varepsilon
  \end{align}
  Let $\varepsilon':= 2\varepsilon$. For all $n>N$
  \begin{align}
    |y_n-L| < \varepsilon
  \end{align}
  Or
  \begin{align}
    \lim_n y_n = L
  \end{align}
\end{proof}

\begin{theorem}
  Positive fundamental sequences are closed over equivalence
  \begin{align}
    \forall x\in\mathcal{F}^+: x\equiv_\mathcal{Z} y \implies y\in\mathcal{F}^+
  \end{align}
\end{theorem}
\begin{proof}
  Since $x$ is in $\mathcal{F}^+$, there is an $r\in\mathcal{Q}^+$ and $N\in\mathbb{Z}^+$ such that for all $n>N$
  \begin{align}
    x_n > r>0
  \end{align}
  Since the rationals are dense, there is a rational $s$ such that
  \begin{align}
    r> s>0\\
    r-s>0
  \end{align}
  Since $x-y\in\mathcal{Z}$, there is an $N'\in\mathcal{Z}^+$ such that
  \begin{align}
    n> N'> N \implies |x_n - y_n| < r-s
  \end{align}
  If $n> N'$
  \begin{align}
    s-r < y_n - x_n
  \end{align}
  Along with $r< x_n$
  \begin{align}
    s- r+ r = s < y_n- x_n +x_n = y_n
  \end{align}
  \begin{align}
    \forall n>N': s>y_n
  \end{align}
  So $y\in\mathcal{F}^+$
\end{proof}

\begin{definition}[The Real Field]
  The \emph{real field} is the partition of equivalence classes in $\mathcal{F}$
  \begin{align}
    \mathbb{R} := \mathcal{F} / \equiv_\mathcal{Z}
  \end{align}
\end{definition}


\begin{definition}[Convergent Sequence]
  A real, convergent sequence converges to $L\in\mathbb{R}$ if
  \begin{align}
    \forall \varepsilon \in \mathbb{R}^+: \exists N \in \mathbb{Z}^+: n > N \implies |x_n - L | < \varepsilon
  \end{align}
  If $L$ exists it is called the \emph{limit} of the function
  \begin{align}
    L = \lim_n x_n
  \end{align}
\end{definition}

\begin{theorem}%[Archimedean Property]
  For every real number there is a greater positive integer.
  \begin{align}
    \forall x\in\mathbb{R}: \exists n \in \mathbb{Z}^+: n^* > x 
  \end{align}
  where $n^*:= m \mapsto n$ and $\mathbb{Z}^+ \subset \mathbb{Q}$
\end{theorem}
\begin{proof}
  Let $x\in\mathbb{R}$ be any real number.
  By definition it is a fundamental sequence of rationals $x\in\mathcal{C}$.
  Every fundamental sequence is bounded, so there exists $b,B\in\mathbb{Q}$, such that, for all $i\in\mathbb{Z}+$
  \begin{align}
    b < x_i < B
  \end{align}
  By the Archimedean property of rationals, there exists an $n\in\mathbb{Z}^+$, such that $n>B$.
  \begin{align}
    x_i < n
  \end{align}
  Consider the constant sequence $n^*: i \mapsto n$
  \begin{align}
    x_i < n = n^*_i
  \end{align}
  Choose any $\varepsilon \in\mathbb{Q}^+$ and any $N\in\mathbb{Z}^+$
  \begin{align}
    i > N \implies x_i < n_i^*
  \end{align}
  which by definition of order on the reals
  \begin{align}
    x< n^*
  \end{align}
\end{proof}

\begin{theorem}
  $\mathbb{Q}^*$ is a dense subset of $\mathbb{R}$
  \begin{align}
    \forall x \in \mathbb{R}, \varepsilon \in \mathbb{Q}^+: \exists r \in \mathbb{Q}^*: |x - r | < \varepsilon
  \end{align}
\end{theorem}
\begin{proof}
  Fix $\varepsilon\in \mathbb{Q}^+$ and $x \in\mathbb{R}$.
  Since $x$ is a rational, fundamental sequence, there is an $N\in\mathbb{Z}^+$ such that
  \begin{align}
    m,n > N \implies |x_m - x_n | < \varepsilon
  \end{align}
  Choose $l>m$, let $r:= n \mapsto x_l$
  \begin{align}
    n> N \implies |x_n - r_n| < \varepsilon
  \end{align}
  By the definition of order on $\mathbb{R}$ and the definition of a limit
  \begin{align}
    \lim_n |x_n - r_n | < \varepsilon
  \end{align}
  By the definition of the absolute value on $\mathbb{R}$
  \begin{align}
    |x- r| < \varepsilon
  \end{align}
\end{proof}

\begin{theorem}
  Every fundamental sequence of rational, constant sequences converges to a fundamental sequence.
  \begin{align}
    \lim_n r^*_n = r
  \end{align}
  where $r^*: n \mapsto ( m \mapsto r_n)$
\end{theorem}
\begin{proof}
  Let $r\in\mathcal{F}$ be a rational, fundamental sequence. Choose some $\varepsilon\in\mathbb{R}^+$, then there is some $N\in\mathbb{Z}^+$ such that
  \begin{align}
    m,n > N \implies |r_m - r_n | <\varepsilon
  \end{align}
  Fixing $m>N$, the definition of a limit gives
  \begin{align}
    \lim_n r_n = r_m
  \end{align}
  \begin{align}
    \left|\lim_n r_n - r_m \right| = 0 < \varepsilon
  \end{align}
  Define $r^*: m  \mapsto (n \mapsto r_m)$, then the limit of the constant sequences $r_m^*$ are the constant
  \begin{align}
    \lim_n r_{mn}^* = r_m
  \end{align}
  Giving
  \begin{align}
    \left|\lim_n r_n - \lim_n r_{mn}^* \right| < \varepsilon
  \end{align}
  % Clean after this
  By the definition of order on $\mathbb{R}$
  \begin{align}
    |r_n^* - r | < \varepsilon
  \end{align}
  Therefor
  \begin{align}
    \forall \varepsilon\in\mathbb{Q}^+: \exists N \in \mathbb{Z}^+: n > N \implies |r_n^* - r| < \varepsilon
  \end{align}
  or more succinctly
  \begin{align}
    \lim_n r_n^* = r
  \end{align}
\end{proof}


\begin{theorem}[Completeness]
  Every real, pinched sequence converges to a real number.
\end{theorem}
\begin{proof}
  Let $(x_n)$ be a real, pinched sequence. Define a sequence of constant, rational sequences $r^*:\mathbb{N}\to\mathbb{Q}^*$ such that
  \begin{align}
    x_n-\frac{1}{n} < r_n^* < x_n+\frac{1}{n}
  \end{align}
  or
  \begin{align}
    |r_n^*-x_n | < \frac{1}{n}
  \end{align}
  For any $m,n\in\mathbb{Z}^+$, by triangle identity
  \begin{align}
    |r^*_m - r^*_n | &= |r^*_m - x_m + x_m - x_n +x_n - r^*_n|\\
    &\leq |r_m^* - x_m| + |x_m - x_n| + |r_n^* - x_n|\\
    &< |x_m - x_n| + \frac{1}{m} + \frac{1}{n}
  \end{align}
  Since $x$ is a fundamental sequence, for every $\varepsilon\in\mathbb{R}^+$ there is an $N\in\mathbb{Z}^+$, such that $m,n>N$
  \begin{align}
    |x_m - x_n | < \frac{\varepsilon}{3}
  \end{align}
  Choose $N'\in\mathbb{Z}^+:N'>\frac{3}{\varepsilon}$, and define $M:=\max\{ N , N'\}$. For all $m, n > M $
  \begin{align}
    m > M &> \frac{3}{\varepsilon} > 0\nonumber\\
    n > M &> \frac{3}{\varepsilon} > 0
  \end{align}
  implies
  \begin{align}
    \frac{1}{m} < \frac{\varepsilon}{3}\nonumber\\
    \frac{1}{n} < \frac{\varepsilon}{3}
  \end{align}
  Applying this to the earlier inequality
  \begin{align}
    |r^*_m - r^*_n| < \frac{\varepsilon}{3} + \frac{\varepsilon}{3} + \frac{\varepsilon}{3} = \varepsilon
  \end{align}
  Hence $r^*$ is a fundamental sequence.
  %Clean after this
  Therefore by the lemma above
  \begin{align}
    \lim_n r^*_n = r
  \end{align}
  where $r^*: n \mapsto ( m \mapsto r_n)$. 
  From the definition of $r^*$, for every $n\in\mathbb{Z}^+$
  \begin{align}
    |r^*_n - x_n | < \frac{1}{n}
  \end{align}
  Since $\mathbb{R}$ is unbounded
  \begin{align}
    \forall \varepsilon \in \mathbb{R}^+: \exists N\in\mathbb{Z}^+ : \frac{1}{N} < \varepsilon
  \end{align}
  If $n>N$, then $\frac{1}{n} < \frac{1}{N} < \varepsilon$
  \begin{align}
    |r^*_n-x_n| <\frac{1}{N} <  \varepsilon
  \end{align}
  Then, for every $\varepsilon\in\mathbb{R}^+$ there is a $N\in\mathbb{Z}^+$ such that if $n>N$
  \begin{align}
    |r^*_n-x_n| < \varepsilon
  \end{align}
  Since $\lim_n r^*_n = r$
  \begin{align}
    \exists N'\in\mathbb{Z}^+: \forall n> N': |r^*_n - r| < \varepsilon
  \end{align}
  Let $M:= \max\{N, N'\}$, if $n>M$, then for all $\varepsilon\in\mathbb{R}^+$
  \begin{align}
    |r^*_n-x_n| < \frac{\varepsilon}{2}\\
    |r^*_n - r| < \frac{\varepsilon}{2}
  \end{align}
  By the triangle identity
  \begin{align}
    |x_n-r | &= |x_n-r^*_n+r^*_n-r|\nonumber\\
    &\leq |r^*_n-x_n| + |r^*_n - r| < \frac{\varepsilon}{2} +  \frac{\varepsilon}{2} = \varepsilon
  \end{align}
  Hence for all $n>M$
  \begin{align}
    |x_n-r| < \varepsilon
  \end{align}
  Or
  \begin{align}
    \lim_n x_n = r
  \end{align}
\end{proof}


\begin{theorem}[Monotone Convergence Theorem]
  A monotonic sequence converges if and only if it is bounded.
\end{theorem}
\begin{proof}
  Let $s$ be an increasing sequence bounded above by $B$.
  \begin{align}
    m>n \implies B\geq s_m \geq s_n
  \end{align}
  Without loss of generality, assume $s$ is non-fundamental
  \begin{align}
    \exists \varepsilon>0, \forall N: \exists m,n> N: |s_m - s_n | > \varepsilon
  \end{align}
  Let $m>n$ (If $m=n$, then $s_m-s_m = 0\ngtr \varepsilon$).
  \begin{align}
    \forall n: \exists m> n: |s_m - s_n | > \varepsilon
  \end{align}
  Since $s$ is increasing
  \begin{align}
    \forall n: \exists m> n \implies s_m - s_n&> \varepsilon\\ 
    s_m &> \varepsilon + s_n
  \end{align}
  Let $p:= m-n \implies m= k+n$
  \begin{align}
    \forall n, \exists p: s_{n+p} > \varepsilon + s_n
  \end{align}
  Define a sequence $k:\mathbb{Z}^+\to \mathbb{Z}^+$
  \begin{align}
    k_1&:=1\\
    k_{i+1}&:= \min\{p | s_{k_i+p} > s_{k_i}+\varepsilon\} + k_i
  \end{align}
  The sequence $s\circ k$ is a subsequence because $k$ is strictly increasing
  \begin{align}
    p\in\mathbb{Z}^+ \implies k_{i+1} > k_i
  \end{align}
  By construction of $k$
  \begin{align}
    s_{k_{i+1}} > s_{k_i} + \varepsilon
  \end{align}
  Let $A:=s_1 = s_{n_1}$.
  Consider the induction hypothesis
  \begin{align}
    \forall i>1: s_{k_i} - A > (i-1) \varepsilon
  \end{align}
  \begin{case}[Base ($i=2$)]
    \begin{align}
      s_{n_{2}} > s_{n_1} + \varepsilon = A + (2-1) \varepsilon
    \end{align}
  \end{case}
  \begin{case}[Induction]
    Assume without loss of generality
    \begin{align}
      s_{k_i} >A + (i-1) \varepsilon
    \end{align}
    Along with
    \begin{align}
      s_{k_{i+1}} - \varepsilon > s_{k_i}
    \end{align}
    Gives
    \begin{align}
      s_{k_{i+1}} - \varepsilon > A+ (i-1) \varepsilon\\
      s_{k_{i+1}} - A > [(i+1)-1] \varepsilon
    \end{align}
  \end{case}
  By induction 
  \begin{align}
    \forall i>1: s_{n_i} - A > (i-1) \varepsilon
  \end{align}
  Choose $N\in\mathbb{Z}^+$ such that
  \begin{align}
    N > \frac{B-A}{\varepsilon} +1
  \end{align}
  By the Archimedean property, $N$ exists.
  \begin{align}
    s_N - A &> \left(\frac{B-A}{\varepsilon} +1 - 1\right) \varepsilon\\
    s_N &> B
  \end{align}
  Which contradicts the assumption that $s$ is bounded above by $B$.
  Therefor if $s$ is bounded above, it is fundamental.
  By the completion theorem, since $s$ is fundamental, it is also convergent.
  Thus it has shown that bounded, increasing sequences are convergent.
  If $s$ is a bounded, monotonic decreasing sequence, then $-s$ is bounded, monotonic increasing sequence, and therefor convergent, which implies $s$ is convergent.
  Thus all bounded, monotonic sequences are convergent.
  The converse is proven by theorem \ref{Thm:ConvergentImpliesBounded}, all convergent sequences are bounded.
\end{proof}

%\begin{theorem}[Nested Interval Theorem]
%\end{theorem}


\section{Summable Sequences}

\begin{theorem}[Squeeze Theorem]
  For all sequences $s$, and all convergent sequences $r, t$ such that
  \begin{align}
    \lim_n r_n = L\nonumber\\
    \lim_n t_n = L
  \end{align}
  and for all $n\in\mathbb{Z}^+$
  \begin{align}
    r_n < s_n < t_n
  \end{align}
  implies that $s$ is a convergent sequence and its limit is given by
  \begin{align}
    \lim_n s_n = L
  \end{align}
\end{theorem}
\begin{proof}
  Fix $\varepsilon> 0 $.
  Then by the definition of the convergent sequence
  \begin{align}
    \exists N_1: n > N_1 \implies |y_n - L | < \varepsilon\nonumber\\
    \exists N_2: n > N_2 \implies |z_n - L| < \varepsilon
  \end{align}
  Let $N:= \max \{ N_1, N_2 \}$, then for all $n> N$
  \begin{align}
    L- \varepsilon < y_n < L + \varepsilon\nonumber\\
    L - \varepsilon < z_n < L + \varepsilon
  \end{align}
  Using the trichotomy property of order
  \begin{align}
    L - \varepsilon < y_n < x_n < z_n < L+\varepsilon
  \end{align}
  or
  \begin{align}
    L - \varepsilon < x_n < L+ \varepsilon
  \end{align}
  Therefor $n>N$ implies
  \begin{align}
    |x_n - L | < \varepsilon
  \end{align}
  or
  \begin{align}
    \lim_n x_n = L
  \end{align}
\end{proof}

\begin{theorem}
  \begin{align}
    \lim_n \frac{1}{n} = 0
  \end{align}
\end{theorem}
\begin{proof}
  Choose an $\varepsilon' \in \mathbb{R}^+$.
  By the Archimedean property
  \begin{align}
    \exists N \in \mathbb{Z}^+: N > \varepsilon '
  \end{align}
  Since both $N>0$ and $\varepsilon'>0$
  \begin{align}
    \frac{1}{N} < \frac{1}{\varepsilon'}
  \end{align}
  Define $\varepsilon:= \frac{1}{\varepsilon'} \in\mathbb{R}^+$
  \begin{align}
    \frac{1}{N} < \varepsilon
  \end{align}
  Let $n>N>0$, then
  \begin{align}
    \frac{1}{N} > \frac{1}{n}
  \end{align}
  By trichotomy
  \begin{align}
    \frac{1}{n} &< \varepsilon\\
    n>N \implies \left| \frac{1}{n} - 0 \right| &< \varepsilon
  \end{align}
  or
  \begin{align}
    \lim_n \frac{1}{n} = 0
  \end{align}
\end{proof}

\begin{lemma}\label{Lemma:ExponentialConvergentSequence}
  If $\{|a_n-L|\}$ converges to zero, then $\{a_i\}$ converges to $L$.
  \begin{align}
    \lim_n |a_n-L| = 0 \implies \lim_n a_n = L
  \end{align}
\end{lemma}
\begin{proof}
  Since $\{|a_n-L|\}$ converges to zero, for any $\varepsilon>0$, there is an $N$
  \begin{align}
    n> N \implies \left| |a_n- L | - 0\right|&< \varepsilon\\
    \left| |a_n - L | \right|&< \varepsilon\\
    n>N\implies |a_n-L| &< \varepsilon
  \end{align}
  Then
  \begin{align}
    \lim_n a_n = L
  \end{align}
\end{proof}


\begin{definition}[Series]
  A \emph{series} of partial sums of $a$ is the sequence
  \begin{align}
    s_n := \sum^n_{i=1} a_i
  \end{align}
\end{definition}

\begin{definition}[Summable Sequence]
  A sequence is \emph{summable} if its series of partial sums converges
  \begin{align}
    \sum^\infty_{i=1} a_i := \lim_n \sum^n_{i=1} a_i
  \end{align}
  The \emph{sum} of a sequence is the limit of its series.
\end{definition}

\begin{theorem}
  Each sequence with the same sign has a monotonic series.
  \begin{itemize}
  \item A non-negative sequence has an increasing series.
  \item A positive sequence has a strictly increasing series.
  \item A non-positive sequence has a decreasing series.
  \item A negative sequence has a strictly decreasing series.
  \end{itemize}
\end{theorem}
\begin{proof}
  Let $\sim$ be the relationship between the sequence $a$ and zero.
  \begin{align}
    \forall n: a_n\sim 0
  \end{align}
  Then %holds only for sums not = 0 
  \begin{align}
    a_{n+1} \sim 0\\
    a_{n+1} + s_n \sim s_n\\
    s_{n+1} \sim s_n
  \end{align}
\end{proof}

\begin{theorem}[Series Linearity]
  If $a$ and $b$ are summable and $c$ and $d$ are some numbers, then $ca+bd$ is summable and
  \begin{align}
    c\sum_{i=1}^\infty a_i + d\sum_{i=1}^\infty b_i = \sum_{i=1}^\infty (ca_i+db_i)
  \end{align}
\end{theorem}
\begin{proof}
  \begin{align}
    c\sum_{i=1}^\infty a_i + d\sum_{i=1}^\infty b_i &= c\lim_{n} \sum^n_{i =1} a_i + d \lim_n \sum^n_{i=1}b_i\\
    &= \lim_n \left(c\sum^n_{i=1} a_i +  d\sum^n_{i=1} b_i \right)\\
    &= \lim_n \sum^n_{i=1} (ca_i+db_i)\\
    &= \sum^\infty_{i=1} (ca_i + db_i)
  \end{align}
\end{proof}

\begin{theorem}[Fundamental Criterion]
  A sequence $a$ is summable if and only if
  \begin{align}
    \forall \varepsilon > 0: \exists N: n, m> N: \left| \sum_{i=m+1}^n a_i \right| < \varepsilon
  \end{align}
\end{theorem}
\begin{proof}
  Fix $\varepsilon>0$. By assumption, there is an $N$ such that
  \begin{align}
    n, m > N \implies \left|\sum^n_{i=m+1} a_i \right| < \varepsilon
  \end{align}
  If $n\leq m$, $\sum^n_{i=m+1} = 0$, which is less than all $\varepsilon>0$.Restrict $n>m$, then
  \begin{align}
    \left| \sum^n_{i=1} a_i - \sum^{m}_{i=1} a_i \right| < \varepsilon
  \end{align}
  which by the completeness theorem, shows that the series of $a$ converges and $a$ is summable.
\end{proof}

\begin{theorem}[Vanishing Condition]
  If $a$ is summable
  \begin{align}
    \lim_n a_n = 0
  \end{align}
\end{theorem}
\begin{proof}
  Since $a$ is summable it has a sum. Let $S$ be the sum of $a$
  \begin{align}
    S:= \sum^\infty_{i=1} a_i
  \end{align}
  Then
  \begin{align}
    \lim_n a_n &= \lim_n \left(\sum^n_{i=1} a_n- \sum^{n-1}_{i=1} a_n \right) \\
    &= \lim_n \sum^n_{i=1} a_n - \lim_n \sum^{n-1}_{i=1} a_n\\
    &= S - S = 0
  \end{align}
\end{proof}

\begin{theorem}[Boundness Criterion]
  A non-negative sequence is summable, if and only if its series is bounded.
\end{theorem}
\begin{proof}
  Let $a$ be a non-negative sequence. Then its series is increasing. By the monotonic convergence theorem, this monotonic series is convergent if and only if it is bounded. Therefor the non-negative sequence is summable if and only if it is bounded.
\end{proof}

\begin{theorem}[Comparison Test]
  If $\forall n:|a_n| \leq b_n$	and $b$ is summable, then $a$ is summable
\end{theorem}
\begin{proof}
  Choose $\varepsilon>0$. By fundamental criterion, there is an $N$, such that $m, n> N$
  \begin{align}
    \left| \sum^n_{i=m} b_i\right| < \varepsilon
  \end{align}
  Since $b$ is non-negative
  \begin{align}
    \sum^n_{i=m} b_i = \left| \sum^n_{i=m} b_i\right| < \varepsilon
  \end{align}
  Using the triangle inequality
  \begin{align}
    \left| \sum_{i=m}^n a_i \right| \leq \sum^n_{i=m} |a_i| \leq \sum_{i=m}^n b_i < \varepsilon
  \end{align}
  By the fundamental criterion, $a$ is then summable.
\end{proof}

\begin{theorem}[Geometric Series]
  If $|a|<1$, then $\langle a^k \rangle$ is summable. The sum is given by
  \begin{align}
    \sum^\infty_{i=0} a^i = \frac{1}{1-a}
  \end{align}
\end{theorem}
\begin{proof}
  Let $\langle s_n \rangle$ be the series of partial sums
  \begin{align}
    s_n := \sum^n_{i=0} a^i
  \end{align}
  
  \begin{align}
    \lim_n s_n = \lim_n \frac{1-a^{n+1}}{1-a} = \frac{1}{1-a}
  \end{align}
\end{proof}

\begin{lemma}
  \begin{align}
    \limsup_n s_n < B \implies \exists N\in\mathbb{N}: \forall n\geq N: s_n \leq B
  \end{align}
\end{lemma}
\begin{proof}
  Let $\langle s_n\rangle$ be a bounded sequence
  \begin{align}
    L : = \limsup_n s_n
  \end{align}
  Define $B$ such that
  \begin{align}
    L < B
  \end{align}
  Suppose that
  \begin{align}
    \forall N\in \mathbb{N}: \exists n > N: s_n > B
  \end{align}
  Define a sequence $k: \mathbb{N} \to \mathbb{N}$, such that
  \begin{align}
    k_1 &= \min \{ n: s_n > B \}\\
    k_{i+1} &= \min \{ n : s_{k_i+n} > B \} + k_i
  \end{align}
  Since $k_{i+1}> k_i$, $\langle s\circ k \rangle$ is a subsequence. By its definition, $\forall i\in \mathbb{N}: s_{k_i} > B > L$.
  Then $\lim_i s_{k_i} > \limsup_n s_n$ which is a contradiction.
  Therefor $\exists N: \forall n \geq N: s_n < B$
\end{proof}


\begin{theorem}[Ratio Test]
  If 
  \begin{align}
    \limsup_n \left|\frac{a_{n+1}}{a_n}\right| < 1
  \end{align}
  Then $a$ is summable.
\end{theorem}
\begin{proof}
  Let $a$ be some sequence such that
  \begin{align}
    \limsup_n \left|\frac{a_{n+1}}{a_n} \right|< 1
  \end{align}
  Define $s$ such that
  \begin{align}
    \limsup_n \left|\frac{a_{n+1}}{a_n} \right|< s< 1
  \end{align}
  By the above lemma, there is an $N$, such that for every $n\geq N$
  \begin{align}
    \left|\frac{a_{n+1}}{a_n}\right| &< s\\
    \left|a_{n+1} \right|&< s |a_n|
  \end{align}
  Induction hypothesis
  \begin{align}
    \forall k \in\mathbb{Z}^+: |a_{N+k} | < s^k | a_N |
  \end{align}
  The base case is given by $n = N$
  \begin{align}
    |a_{N+1} | < s | a_N|
  \end{align}
  Assume without loss of generality $|a_{N+k}| < s^k |a_N|$, multiplying by $s>0$
  \begin{align}
    |a_{N+k} | s < s^{k+1} |a_n|
  \end{align}
  Applying the inequality from before
  \begin{align}
    |a_{N+k+1} | < |a_{N+k} | s < s^{k+1} |a_n|
  \end{align}
  Thus, by induction
  \begin{align}
    \forall k \in\mathbb{Z}^+: |a_{N+k} | < s^k | a_N |
  \end{align}
  By the comparison test, $a_{N+k}$ is summable, therefor $a_n$ is summable.
\end{proof}

\begin{definition}[Absolutely Summable]
  A sequence is \emph{absolutely summable} if $\langle |a_n| \rangle$ is summable.
\end{definition}


\begin{theorem}
  An absolutely summable sequence is summable.
\end{theorem}
\begin{proof}
  By assumption the series
  \begin{align}
    \sum^\infty_{i=1} |a_i|
  \end{align}
  converges. Then by the fundamental criterion, for every $\varepsilon>0$, there is an $N$, such that
  \begin{align}
    m,n > N \implies \left| \sum^n_{i=m+1} |a_i| \right|&< \varepsilon\\
    \sum^n_{i=m+1} |a_i| \leq \left| \sum^n_{i=m+1} |a_i| \right|&< \varepsilon
  \end{align}
  By the triangle identity
  \begin{align}
    \left|\sum^n_{i=m+1} a_i\right|\leq \sum^n_{i=m+1} |a_i|< \varepsilon
  \end{align}
  Then, for every $\varepsilon>0$, there is an $N$, such that
  \begin{align}
    m,n >N \implies \left|\sum^n_{i=m+1} a_i\right|
  \end{align}
  By the fundamental criterion, $a$ is summable.
\end{proof}

\begin{lemma}
  If $\langle a_n \rangle$ is non-negative and decreasing, then
  \begin{align}
    \forall m,n\in\mathbb{N}: \left| \sum_{i=m}^n  (-1)^{i} a_i \right| \leq a_m
  \end{align}
\end{lemma}
\begin{proof}
  Fix $m\in\mathbb{N}$. Let the induction hypothesis be that for all $n\in\mathbb{N}$
  \begin{align}
    0 \leq \sum^n_{i=m} (-1)^{i-m}a_i \leq a_m
  \end{align}
  \begin{case}[Basis for Induction $P(m)$]
    Since $\langle a_n \rangle$ is non-negative
    \begin{align}
      0 \leq a_m = (-1)^{m-m} a_m = \sum^m_{i=m} (-1)^{i-m} a_i
    \end{align}
  \end{case}
  
  \begin{case}[Basis for Induction $P(m+1)$]
    Since $\langle a_n \rangle$ is non-negative
    \begin{align*}
      a_m &\geq -a_{m+1} + a_m \geq 0\\
      a_m &\geq (-1)^{m+1-m} a_{m+1} + (-1)^{m-m} a_m \geq 0\\
      a_m &\geq \sum_{i=m}^{m+1} (-1)^{i-m} a_i \geq 0
    \end{align*}
  \end{case}
  For the induction step, suppose for $\forall k: m\leq k \leq n$
  \begin{align}
    0 \leq \sum^k_{i=m} (-1)^{i-m} a_i \leq a_m
  \end{align}
  \begin{case}
    Further restrict $n-m$ to odd integers. Since $\langle a_n \rangle$ is non-negative $a_{n+1} \geq 0$ along with the induction hypothesis
    \begin{align}
      0 &\leq a_{n+1} + \sum^n_{i=m} (-1)^{i-m} a_i\\
      &= (-1)^{n+1-m} a_{n+1} +  \sum^n_{i=m} (-1)^{i-m} a_i\\
      &= \sum^{n+1}_{i=m} (-1)^{i-m} a_i
    \end{align}
    Since $\langle a_n\rangle$ is decreasing
    \begin{align}
      a_{n+1} \leq a_n \implies a_{n+1} - a_n \leq 0
    \end{align}
    Along with the induction hypothesis
    \begin{align}
      a_m &\geq a_{n+1} - a_n + \sum^{n-1}_{i=m} (-1)^{i-m} a_i\\
      &= (-1)^{n+1-m} a_{n+1} + (-1)^{n-m} a_n + \sum_{i=m}^{n-1} (-1)^{i-m} a_i\\
      &= \sum^{n+1}_{i=m} (-1)^{i-m} a_i
    \end{align}
  \end{case}
  \begin{case}
    Now restrict $n-m$ to only even integers.
    Since $\langle a_n \rangle$ is decreasing
    \begin{align}
      a_n \geq a_{n+1} \implies a_n - a_{n+1} \geq 0
    \end{align}
    Applying the induction hypothesis
    \begin{align}
      0 &\leq -a_{n+1} + a_n + \sum_{i=m}^{-1} (-1)^{i-m} a_i\\
      &=(-1)^{n+1-m} a_{n+1} + (-1)^{n-m} a_n + \sum_{i=m}^{n-1} (-1)^{i-m} a_i\\
      &= \sum^{n+1}_{i=m} (-1)^{i-m} a_i
    \end{align}
    Since $\langle a_n \rangle$ is non-negative $-a_{n+1} \leq 0$. Applying the hypothesis gives
    \begin{align}
      a_m &\geq -a_{n+1} + \sum^n_{i=m} (-1)^{i-m} a_i \\
      &=(-1)^{n+1-m} a_{n+1} + \sum^n_{i=m} (-1)^{i-m}a_i \\
      &=\sum^{n+1}_{i=m} (-1)^{i-m} a_i
    \end{align}
  \end{case}
  Therefor for all $n$, $P(n)\implies P(n+1)$. By induction, for all $n$,
  \begin{align}
    0 \leq \sum^n_{i=m} (-1)^{i-m}a_i \leq a_m
  \end{align}
  Since the sum is greater than zero
  \begin{align}
    \sum^n_{i=m} (-1)^{i-m}a_i &= \left|\sum^n_{i=m} (-1)^{i-m}a_i \right|\\
    &=\left|(-1)^m \sum^n_{i=m} (-1)^{i-m}a_i \right|\\
    &=\left| \sum^n_{i=m} (-1)^{i}a_i \right| \leq a_m
  \end{align}
\end{proof}

\begin{theorem}[Alternating Series Test]
  If $a$ is decreasing, non-negative and converges to zero, then $\left\langle (-1)^n a_n \right\rangle$ is summable.
\end{theorem}
\begin{proof}
  Let $\langle s_n \rangle$ be a series of partial sums
  \begin{align}
    s_n:= \sum^n_{i=0} (-1)^n a_n
  \end{align}
  Fix $\varepsilon>0$. Since $\lim_n a_n =0$, choose some $N$ such that if $n>N$
  \begin{align}
    |a_n| < \varepsilon
  \end{align}
  Choose $n$ and $m$ such that $n>m>N$, then
  \begin{align}
    %|s_n-s_m| &= \left| \sum_{i=0}^n (-1)^i a_i - \sum_{i=0}^m (-1)^i a_i \right|\\
    %&=
    \left| \sum_{i=m}^n (-1)^i a_i \right| \leq a_{m} \leq |a_{m} | < \varepsilon
  \end{align}
  By the fundamental criterion, $\langle a_n \rangle$ is summable.
\end{proof}


\section{Functional Limits}

\begin{definition}[Limit]
  The \emph{limit} $L$ of a function $f: [a,b]\to \mathbb{R}$ at $c$ is
  \begin{align}
    \forall \varepsilon> 0: \exists \delta >0: \forall x\in[A,B]:  0<|x-c| < \delta \implies |f(x)- L| < \varepsilon
  \end{align}
  The limit at $c$ of $f$ is denoted by
  \begin{align}
    \lim_{x\to c} f(x) = L 
  \end{align}
\end{definition}

\begin{theorem}[Sequential Limit Criterion]
  \begin{align}
    \lim_{x\to c} f(x) = L
  \end{align}
  if and only if, for all sequences such that $\forall i: x_i \neq c$:
  \begin{align}
    \lim_n x_n = c \implies \lim_n f(x_n) = L
  \end{align}
\end{theorem}
\begin{proof}
  Choose a sequence $x$ such that $\forall x_n\neq c$ and
  \begin{align}
    \lim_n x_n = c
  \end{align}
  Fix $\varepsilon > 0$. By assumption $\exists \delta>0$ such that for every $x$
  \begin{align}
    0 < |x-c | < \delta \implies |f(x) - L| < \varepsilon
  \end{align}
  Since $x_n$ converges to $c$, there is an $N$ such that
  \begin{align}
    n> N \implies |x_n - c| < \delta
  \end{align}
  Using the definition of the functional limit, this implies
  \begin{align}
    |f(x_n) - L | < \varepsilon
  \end{align}
  Thus
  \begin{align}
    \lim_n f(x_n) = L
  \end{align}
  For the converse, assume
  \begin{align}
    \lim_{x\to c} f(x) \neq L
  \end{align}
  Then there is an $\varepsilon >0$ such that for every $\delta>0$ there is an $x$ such that
  \begin{align}
    0< |x-c| < \delta \land |f(x) - L| \geq \varepsilon
  \end{align}
  Choose such an $\varepsilon$. Generate a sequence such that
  \begin{align}
    0< |x_n-c| < \frac{1}{n} \land |f(x_n) - L| \geq \varepsilon
  \end{align}
  By the squeeze theorem
  \begin{align}
    \lim_n x_n = c
  \end{align}
  The other part of the generator
  \begin{align}
    \forall n: |f(x_n) - L | \geq \varepsilon
  \end{align}
  implies
  \begin{align}
    \lim_n f(x_n) \neq L
  \end{align}
  Thus
  \begin{align}
    \lim_n x_n = c \implies \lim_n f(x_n) =L \implies \lim_{x\to c} f(x) =L
  \end{align}
\end{proof}

\begin{theorem}[Squeeze Theorem]
  If
  \begin{align}
    \lim_{x\to c} f(x) = L\\
    \lim_{x\to c} g(x) = L
  \end{align}
  and
  \begin{align}
    f(x) \leq h(x) \leq g(x)
  \end{align}
  Then
  \begin{align}
    \lim_{x\to c} h(x) = L
  \end{align}
\end{theorem}

\begin{theorem}[Uniqueness]
  If a function converges to a limit at some point, then that limit is unique.
\end{theorem}

\begin{theorem}
  \begin{align}
    \lim_{x\to c} f(x) = F\\
    \lim_{x\to c} g(x) = G
  \end{align}
  implies
  \begin{align}
    \lim_{x\to c} \left[f(x) + g(x) \right] = F+G 
  \end{align}
\end{theorem}

\begin{theorem}
  \begin{align}
    \lim_{x\to c} f(x) = F\\
    \lim_{x\to c} g(x) = G
  \end{align}
  implies
  \begin{align}
    \lim_{x\to c} f(x)g(x) = FG
  \end{align}
\end{theorem}

\section{Continuous Functions}


\begin{definition}[Continuous Function]
  A function $f$ is continuous at $a$ if
  \begin{align}
    \lim_{x\to a} f(x) = f(a)
  \end{align}
\end{definition}

\begin{theorem}[Sequential Limit Criterion]
  For continuous functions $f$
  \begin{align}
    \lim_n f\left(x_n\right) = f\left(\lim_n x_n\right)
  \end{align}
\end{theorem}

\begin{theorem}[$\varepsilon$-$\delta$ Condition]
\end{theorem}

\begin{theorem}[Closed Set Condition]
  If $f:M\to N$ is continuous and $M$ is closed, then the pre-image of each closed set in $N$ is closed in $M$.
\end{theorem}
\begin{proof}
  Let $I$ be a closed subset of $N$. Define a sequence $s:\mathbb{N}\to I\subset M$ such that
  \begin{align}
    \lim_n s_n = L \in M
  \end{align}
  By sequential convergence condition
  \begin{align}
    \lim_n f(s_n) = f(\lim_n s_n ) = f(L)
  \end{align}
  Since the set $I$ is closed
  \begin{align}
    \lim_n f(s_n) \in I \implies f(L) \in I
  \end{align}
  So $L \in f^{-1} [I]$ which implies that $f^{-1} I$ is closed
\end{proof}

%\begin{corollary}[Open Set Condition]
%If $f:M\to N$ is continuous, then the pre-image of each open set in $N$ is open in $M$.
%\end{corollary}
%\begin{proof}
%Let $I$ be an open subset of $N$, then $N-I$ is closed. By the closed set condition $f^{-1}[N-I]$ is closed. So $M-f^{-1} [N-I]$ is open.
%\end{proof}

\begin{theorem}
  If $f:M\to N$ is continuous and $I$ is compact in $M$, then the image of $f$ restricted to $I$ is compact
\end{theorem}
\begin{proof}
  Let $b:\mathbb{N}\to f[I]$ and define $a:\mathbb{N}\to I$ such that
  \begin{align}
    f(a_n) = b_n
  \end{align}
  Since $I$ is compact, $a$ has a subsequential limit in $I$
  \begin{align}
    L:= \lim_i a_{n_i} \in I
  \end{align}
  Since $f$ is continuous
  \begin{align}
    \lim_i b_{n_i} = \lim_i f(a_{n_i}) = f(L) \in f[I]
  \end{align}
  There for the image of $f$ is restricted to $A$ is compact.
\end{proof}


\begin{lemma}\label{Lemma:IntermediateValueTheorem}
  \begin{align}
    a_{n+1} := \frac{a_n}{2} \implies a_n = \frac{1}{2^{n-1}} a_1
  \end{align}
\end{lemma}
\begin{proof}
  The induction hypothesis is that, $\forall n\in\mathbb{Z}^+$
  \begin{align}
    a_n = \frac{1}{2^{n-1}} a_1
  \end{align}
  \begin{case}[Base]
    The base case is trivially true
    \begin{align}
      a_1 &= \frac{1}{2^{1-1}} a_1
    \end{align}
  \end{case}
  \begin{case}[Induction]
    Assume without loss of generality
    \begin{align}
      a_n &= \frac{1}{2^{n-1}} a_1
    \end{align}
    From the opening assumption
    \begin{align}
      a_{n+1} &= \frac{a_n}{2}\\
      &=\frac{1}{2} \frac{1}{2^{n-1}} a_1\\
      &= \frac{1}{2^{n+1-1}} a_1
    \end{align}
  \end{case}
  By induction, $\forall n \in \mathbb{Z}^+$
  \begin{align}
    a_n = \frac{1}{2^{n-1}} a_1
  \end{align}
\end{proof}

\begin{theorem}[Intermediate Value Theorem]
  If $f:[A,B] \to \mathbb{R}$ is continuous and $f(A)\leq u \leq f(B)$ or $f(A)\geq u \geq f(B)$, then there exists $C\in[A,B]: f(C)= u$
\end{theorem}
\begin{proof}
  Define four sequences, $a$, $b$, $c$, and $e$ with the following algorithm
  \begin{align}
    a_1&:=A\\
    b_1&:=B\\
    c_n &:= \frac{1}{2} (a_n + b_n)
    e_n = b_n - a_n
  \end{align}
  If $f(c_n) < u$:
  \begin{align}
    a_{n+1} &= c_n\\
    b_{n+1} &= b_n
  \end{align}
  Else:
  \begin{align}
    a_{n+1} &= a_n\\
    b_{n+1} &= c_n
  \end{align}
  Using the algorithm, if $f(c_n) = u$, then there  is a $C=c_n\in[A,B]: f(C) = u$, and we are done.
  For the rest of this proof, we will assume that $\forall n: f(c_n) \neq u$.
  Finishing off the algorithm, if $f(c_n) > u$
  \begin{align}
    a_{n+1} &= a_n\\
    b_{n+1} &= c_n = \frac{a_n+b_n}{2}\\
    e_{n+1} &= \frac{a_n+b_n}{2} - a_n= \frac{b_n- a_n}{2}
  \end{align}
  If $f(c_n) < u$
  \begin{align}
    a_{n+1} &= c_n = \frac{a_n+b_n}{2} \\
    b_{n+1} &= b_n\\
    e_{n+1} &= b_n - \frac{a_n + b_n}{2} = \frac{b_n - a_n}{2}
  \end{align}
  Since $\forall n:a_n<c_n<b_n$, $a_n$ is increasing and $b_n$ is decreasing. Since they are also bounded by $(A,B)$, by the monotone convergence theorem, they both converge.
  From the definition of $e_n$
  \begin{align}
    e_n = b_n- a_n
  \end{align}
  So
  \begin{align}
    e_{n+1} = \frac{e_n}{2}
  \end{align}
  By the lemma \ref{Lemma:IntermediateValueTheorem}
  \begin{align}
    e_n = \frac{1}{2^{n-1}} e_1 = \frac{1}{2^{n-1}} \left(B-A\right)		
  \end{align}
  which converges to zero. 
  Then $\lim_n b_n = \lim_n a_n$. Let $C:=\lim_n a_n$, for continuous functions $f$
  \begin{align}
    \lim_n f(a_n) = f(C) = \lim_n f(b_n)
  \end{align}
  By the squeeze theorem, since $\forall n: f(a_n) \leq u \leq f(b_n)$
  \begin{align}
    f(C) = u
  \end{align}
  Thus it is proved that there exists such a $C$ for $f(A) \leq u \leq f(B)$.
  If $f(A) \geq u \geq f(B)$, then 
  \begin{align}
    - f(A) \leq - u \leq -f(B)
  \end{align}
  Let $g(x) = -f(x)$ and $v= -u$, then
  \begin{align}
    g(A) \leq v \leq g(B)
  \end{align}
  By the earlier result, $\exists C\in [A,B]: g(C) = v$, then $-f(C) = -u$ or $f(C) = u$.
\end{proof}

\begin{theorem}
  If $f:M \to N$ is continuous and $S$ is compact in $M$, then $f[S]$ is compact.
\end{theorem}

\begin{corollary}[Bounded Value Theorem]
  If $f:M \to N$ is continuous and $S$ is compact in $M$, then $f[S]$ is bounded.
\end{corollary}

\begin{theorem}[Bounded Value Theorem]
  If $f:[a,b]\to \mathbb{R}$ is a continuous function, then $f$ is bounded.
  \begin{align}
    \exists C: \forall x\in [a,b]: m\leq f(x) \leq M
  \end{align}
\end{theorem}
\begin{proof}
  Let $f:[A,B]\to \mathbb{R}$ be any continuous function.
  Assume without loss of generality
  \begin{align}
    \forall C: \exists x\in{a,B}: f(x) > C
  \end{align}
  Since $\mathbb{Z}^+\subset \mathbb{R}$
  \begin{align}
    \forall n\in\mathbb{Z}^+, \exists x \in [A,B]: f(x) >n
  \end{align}
  Then there is a sequence $x$ such that for all $n$
  \begin{align}
    x_n\in[A,B] \land f(x_n)>n
  \end{align}
  By the point peak theorem, there is a monotonic subsequence $x\circ k$.
  Since $x$ is bounded, the subsequence $x\circ k$ is bounded.
  Since $x\circ k$ is bounded and monotone $x\circ k$ converges by the monotone convergence theorem.
  Since $f$ is continuous, $f(x\circ k)$ also converges.
  Since $x\circ k$ is a subsequence
  \begin{align}
    k_i \geq i
  \end{align}
  By trichotomy, for all $i$
  \begin{align}
    f(x_{k_i}) > i
  \end{align}
  Thus $f(x\circ k)$ diverges, which is a contradiction. 
  Therefore $f$ is bounded above.
\end{proof}

\begin{theorem}[Extreme Value Theorem]
  If $f:[A,B]\to \mathbb{R}$ is continuous
  \begin{align}
    \exists M\in[A,B], \forall x\in[A,B]: f(x)\leq f(M)
  \end{align}
\end{theorem}
\begin{proof}
  By the boundedness theorem, there is a bound, $B$ to the continuous function. Define four sequences $a$, $b$, $c$, $e$ such that
  \begin{align}
    a_1&:= f(A)\\
    b_1&:=B
  \end{align}
  For all $n$:
  \begin{align}
    c_n:= \frac{a_n+b_n}{2}\\
    e_n := b_n - a_n
  \end{align}
  If $c_n$ is a bound of the continuous function ($\forall x\in[A,B]: c_n\geq f(x)$):
  \begin{align}
    a_{n+1} := a_n\\
    b_{n+1} := c_n
  \end{align}
  Else
  \begin{align}
    a_{n+1} := c_n\\
    b_{n+1} := b_n
  \end{align}
  If $c_n$ is a bound
  \begin{align}
    a_{n+1} &= a_n\\
    b_{n+1} &= c_n = \frac{a_n+b_n}{2}\\
    e_{n+1} &= \frac{a_n+b_n}{2} - a_n= \frac{b_n- a_n}{2}
  \end{align}
  If $c_n$ is not a bound
  \begin{align}
    a_{n+1} &= c_n = \frac{a_n+b_n}{2} \\
    b_{n+1} &= b_n\\
    e_{n+1} &= b_n - \frac{a_n + b_n}{2} = \frac{b_n - a_n}{2}
  \end{align}
  Since $\forall n:a_n<c_n<b_n$, $a_n$ is increasing and $b_n$ is decreasing. Since they are also bounded by $(A,B)$, by the monotone convergence theorem, they both converge.
  From the definition of $e_n$
  \begin{align}
    e_n = b_n- a_n
  \end{align}
  So
  \begin{align}
    e_{n+1} = \frac{e_n}{2}
  \end{align}
  By the lemma \ref{Lemma:IntermediateValueTheorem}
  \begin{align}
    e_n = \frac{1}{2^{n-1}} e_1 = \frac{1}{2^{n-1}} \left(B-A\right)		
  \end{align}
  which converges to zero. Therefor
  \begin{align}
    \lim_n (b_n - a_n) &= 0\\
    L:=\lim_n b_n &= \lim_n a_n
  \end{align}
\end{proof}

\section{Connected Spaces}

\begin{definition}[Connected Space]
  A set is \emph{connected}, if it has no non-empty, proper, clopen subset. A set that is not connected is \emph{disconnected}.
\end{definition}

\begin{theorem}
  If $f:M\to N$ is continuous, and $M$ is connected, then $f[M]$ is connected.
\end{theorem}
\begin{proof}
  Suppose $f[M]$ is disconnected. Then there is some non-empty, proper, clopen subset $A$.
  \begin{align}
    A \subsetneq f[M]
  \end{align}
  The pre-image of $A$, then is a non-empty, proper subset of $M$
  \begin{align}
    f^{-1}[A] \subsetneq M
  \end{align}
  By the open and closed set conditions, the preimage $f^{-1}[A]$ is clopen.
  But by assumption, there is no non-empty, proper, clopen subset. So $f[M]$ must be connected.
\end{proof}

\begin{lemma}
  If $I$ is a connected interval on $\mathbb{R}$,
  \begin{align}
    \forall a,b\in I: a< b \implies \exists c\in I: a<c<b
  \end{align}
\end{lemma}
\begin{proof}
  Suppose its not,
  \begin{align}
    \exists a,b,\in I: a<b \land c\in (-\infty, a] \cup [b, \infty)
  \end{align}
  Let $J:= (-\infty, a] \cup [b, \infty)$. The subsets of $J$, $(-\infty, a]$ and its conjugate $[b,\infty)$ are both closed, therefor both are clopen, and $I$ is disconnected.
\end{proof}

\begin{theorem}[Generalized Intermediate Value Theorem]
  If $f:M\to \mathbb{R}$ is continuous and $I \subset M$ is connected, then
  \begin{align}
    \forall a,b\in I: &f(a) < f(b) \nonumber\\
    \implies \exists c\in I:& f(a) < f(c) < f(b)
  \end{align}
\end{theorem}
\begin{proof}
  Since $I$ is connected, $f[I]$ is connected. Hence there is a $\gamma\in f[I]$
  \begin{align}
    f(a) < \gamma< f(b)
  \end{align}
  Therefor there is a $c\in I$ such that
  \begin{align}
    \exists c \in I : f(a) < f(c) < f(b)
  \end{align}
\end{proof}

\section{Differentiable Functions}

\begin{definition}[Differentiable Function]
  A function $f$ is \emph{differentiable} at $a$ if the following exists.
  \begin{align}
    f'(a) = \lim_{x\to a} \frac{f(x) - f(a)}{x-a}
  \end{align}
  This limit $f'(a)$ is called the \emph{derivative} of $f$ at $a$
\end{definition}

\begin{theorem}
  Every differentiable function is continuous.
\end{theorem}
\begin{proof}
  \begin{align}
    f(x) - f(a) &= \frac{f(x)-f(a)}{x-a} (x-a)\\
    \lim_{x\to a} \left[f(x) - f(a)\right] &= \lim_{x\to a} \left[\frac{f(x) -f(a)}{x-a} (x-a)\right]
  \end{align}
  If $f:[A,B] \to \mathbb{R}$ is differentiable at $a\in [A,B]$, then $f'(a)$ exists.
  \begin{align}
    f'(a) = \lim_{x\to a} \frac{f(x)-f(a)}{x-a}
  \end{align}
  \begin{align}
    \lim_{x\to a} \left[f(x)- f(a)\right] &= f'(a) \cdot 0\\
    &= 0\nonumber\\
    \lim_{x\to a} f(x) &= \lim_{x\to a} f(a)\nonumber\\
    \lim_{x\to a} f(x) &= f(a)
  \end{align}
  Thus $f(x)$ is continuous at $a$
\end{proof}

\begin{theorem}[Chain Rule]
  If $f:[a,b] \to I$ and $g:I \to \mathbb{R}$ are differentiable, then
  \begin{align}
    (f\circ g)'(x) = (f'\circ g)(x) g'(x)
  \end{align}
\end{theorem}
\begin{proof}
  Define $u,v$ such that
  \begin{align}
    u(t):= f'\circ g(x) - \frac{(f\circ g)(x) - (f\circ g)(t)}{g(x)-g(t)}\\
    v(t):= g'(x) - \frac{g(x)-g(t)}{x-t}
  \end{align}
  The limit of $u$ at $x$ is then
  \begin{align}
    \lim_{t\to x} u(t) &= (f'\circ g)(x) - \lim_{t\to x} \frac{(f\circ g)(x) - (f\circ g)(t)}{g(x)-g(t)}\\
    &= (f'\circ g)(x) - \lim_{g(t)\to g(x)} \frac{(f\circ g)(x) - (f\circ g)(t)}{g(x)-g(t)}\\
    &= (f'\circ g)(x) - (f'\circ g) (x) =0
  \end{align}
  And the limit of $v$ at $x$ is also
  \begin{align}
    \lim_{t\to x} v(t) &=g'(x) - \lim_{t\to x} \frac{g(x)-g(t)}{x-t} \\
    &= g'(x) - g'(x) = 0
  \end{align}
  Rearranging the definition of $v$ gives
  \begin{align}
    g(x) - g(t) = (x-t) \left[g'(x) - v(t)\right]
  \end{align}
  and the definition of $u$
  \begin{align}
    (f\circ g) (x) - (f\circ g) (t) = \left[g(x) - g(t)\right] \left[f'\circ g(x) - u(t)\right]
  \end{align}
  Substituting gives
  \begin{align}
    (f\circ g) (x) - (f\circ g) (t) = (x-t) \left[g'(x) - v(t)\right] \left[f'\circ g(x) - u(t)\right]
  \end{align}
  The definition of the limit of the composite function gives
  \begin{align}
    (f\circ g)'(x) &= \lim_{t\to x} \frac{(f \circ g)(x) - (f\circ g) (t) }{x-t}\\
    &= \lim_{t \to x} \frac{x-t}{x-t}  \left[g'(x) - v(t)\right] \left[f'\circ g(x) - u(t)\right]\\
    &= (f'\circ g) (x) \cdot g'(x)
  \end{align}
\end{proof}

\begin{lemma}%[Rolle's Theorem]
  If $f$ is continuous on $[a,b]$ and differentiable on $(a,b)$, then
  \begin{align}
    f(a) = f(b) \implies \exists c\in(a,b) : f'(c) = 0
  \end{align}
\end{lemma}
\begin{proof}
  Since $f$ is continuous on $[a,b]$, $f$ achieves a maximum $M\in[a,b]$ and a minimum $m\in[a,b]$.
  \begin{case}
    Suppose $m$ and $M$ are both end points of $[a,b]$. Then 
    \begin{align}
      f(a) = f(b) \implies f(m) = f(M)
    \end{align}
    Since $f(m)$ and $f(M)$ are minimum and maximums of $f$ respectively
    \begin{align}
      \forall x\in[a,b]: f(m)\leq f(x) \leq f(M)
    \end{align}
    along with $f(m)=f(M)$, implies
    \begin{align}
      \forall x\in[a,b]: f(m) = f(x) = f(M)
    \end{align}
    So $f$ is a constant function on $[a,b]$, and the derivative of the constant function is
    \begin{align}
      \forall x\in(a,b): f'(x) = 0
    \end{align}
  \end{case}
  \begin{case}
    Now suppose that $m$ or $M$ is not an end point of $[a,b]$.
    Then there is a $c\in(a,b)$ such that $f$ has a local extrema at $c$. Since the derivative at a local extrema is zero.
    \begin{align}
      f'(c) = 0
    \end{align}
  \end{case}
  So for all cases,
  \begin{align}
    \exists c\in(a,b) : f'(c) = 0
  \end{align}
\end{proof}


\begin{theorem}[Generalized Mean Value Theorem]
  If $f$ and $g$ are continuous on $[a,b]$ and differentiable on $(a,b)$, then
  \begin{align}
    \exists c\in (a,b) : \left[f(b)-f(a)\right] g'(c) = \left[g(b)-g(a)\right] f'(c)
  \end{align}
\end{theorem}
\begin{proof}
  Define $h$ such that
  \begin{align}
    h(x):= \left[f(b)-f(a)\right] g(x) - \left[ g(b) - g(a) \right] f(x)
  \end{align}
  Since $f$ and $g$ both are continuous on $[a,b]$ and differentiable on $(a,b)$, $h$ is continuous on $[a,b]$ and differentiable on $(a,b)$. So the derivative of $h$ on $(a,b)$ is
  \begin{align}
    h'(x) = \left[f(b)-f(a)\right] g'(x) - \left[ g(b) - g(a) \right] f'(x)
  \end{align}
  At the end points
  \begin{align}
    h(a) &= f(b) g(a) - f(a) g(a) - g(b) f(a) +g(a)f(a)\nonumber\\
    &=f(b)g(a) - f(a) g(b)\nonumber
  \end{align}
  And
  \begin{align}
    h(b) &= f(b) g(b) - f(a) g(b) - g(b) f(b) + g(a) f(b)\nonumber\\
    &=f(b)g(a) - f(a) g(b)
  \end{align}
  So
  \begin{align}
    h(a)=h(b)
  \end{align}
  By the above lemma
  \begin{align}
    \exists c\in(a,b) : h'(c) =0
  \end{align}
  So there is a $c\in(a,b)$ such that
  \begin{align}
    h'(c) = 0 = \left[f(b)-f(a)\right] g'(c) - \left[ g(b) - g(a) \right] f'(c)
  \end{align}
  or
  \begin{align}
    \left[f(b)-f(a)\right] g'(c) = \left[ g(b) - g(a) \right] f'(c)
  \end{align}
\end{proof}

\begin{theorem}[L'h\^{o}pital's Rule]
  If $f$ and $g$ are differentiable functions and
  \begin{align}
    \lim_{x\to a} f(x) &= 0\nonumber\\
    \lim_{x\to a} g(x) &= 0
  \end{align}
  Then
  \begin{align}
    \lim_{x\to a} \frac{f(x)}{g(x)}= \lim_{x\to a} \frac{f'(x)}{g'(x)} \label{Equation:LHopitalRatio}
  \end{align}
\end{theorem}
\begin{proof}
  Choose some $\varepsilon$ and $\varepsilon'$ such that $\varepsilon>\varepsilon'>0$. By \ref{Equation:LHopitalRatio}, there is some $\delta>0$ such that $\forall x\in[a,b]$
  \begin{align}
    0<|x-c|< \delta \implies \left|\frac{f'(x)}{g'(x)} - L\right|<\varepsilon'
  \end{align}
  For all $(x,y) \subset (c,c+\delta)$ or $(x,y)\subset (c-\delta, c)$, by generalized mean value theorem, there is a $t\in(x,y)$ such that
  \begin{align}
    \frac{f(x)-f(y)}{g(x)-g(y)} = \frac{f'(t)}{g'(t)}
  \end{align}
  Hence 
  \begin{align}
    \left|\frac{f(x)-f(y)}{g(x)-g(y)} - L\right|< \varepsilon'\nonumber\\
    L - \varepsilon'< \frac{f(x)-f(y)}{g(x)-g(y)} < L +  \varepsilon'\nonumber\\
    L-\varepsilon'\leq \lim_{x\to c} \frac{f(x) - f(y)}{g(x)-g(y)}\leq L + \varepsilon'
  \end{align}
  Assume that $\lim_{x\to c} f(x) = 0$ and $\lim_{x\to c} g(x) = 0$, then
  \begin{align}
    L - \varepsilon' \leq \frac{f(y)}{g(y)} \leq L + \varepsilon'\\
    \left|\frac{f(y)}{g(y)} - L\right|\leq \varepsilon'< \varepsilon
  \end{align}
  So for every $\varepsilon>0$, there is a $\delta>0$ such that for all $y\in[a,b]$
  \begin{align}
    |y-c| < \delta \implies \left|\frac{f(y)}{g(y)} - L\right|< \varepsilon
  \end{align}
  Or
  \begin{align}
    \lim_{y\to c} \frac{f(x)}{g(x)} = L
  \end{align}
\end{proof}

\section{Integrable Functions}

\begin{definition}[Partition]
  A \emph{partition} on $[a,b]$ is a finite, increasing sequence $x:[0,n]\cap\mathbb{N} \to [a,b]$, such that $x_0 = a$ and $x_n = b$.
  The sequence of lengths of partitions $\Delta x_i:[1,n]\cap\mathbb{N}\to\mathbb{R}$ is defined simply as
  \begin{align*}
    \Delta x_i := x_i - x_{i-1}.
  \end{align*}
  The \emph{mesh} of a partition is given by
  \begin{align*}
    \mesh x := \max_i \Delta x_i.
  \end{align*}
  A \emph{partition pair} is a tuple of sequences $(x, t)$, where $x$ is a partition and $t:[0,n]\cap\mathbb{N} \to [a,b]$ is a finite sequence of sample points, that satisfy $x_{i-1}\leq t_i \leq x_i$, for all $i\in[0,n]\cap\mathbb{N}$.  
\end{definition}

\begin{definition}[Partition Refinement]
  The sequence $x^*:[0,n^*]\cap\mathbb{N} \to [a, b]$ is a \emph{refinement} of the partition $x:[0, n]\cap\mathbb{N} \to [a, b]$ if every point of $x$ appears in $x^*$.
  That is, there exists a sequence $p: [0, n]\cap\mathbb{N} \to [0, n^*]\cap\mathbb{N}$, such that
  \begin{align*}
    x = x^* \circ p.
  \end{align*}
\end{definition}

\begin{definition}[Partition Sum]
  The \emph{partition sum} of a function on $f$ with partition pair $(x, t)$ is given by
  \begin{align*}
    R(f,x,t):= \sum_{i=1}^n f(t_i) \Delta x_i.
  \end{align*}
\end{definition}

\begin{definition}[Integral]
  A function is \emph{integrable} on $[a,b]$ if there is an $I\in\mathbb{R}$ such that, for all $\varepsilon>0$ there is a $\delta>0$ such that for all partition pairs $(x, t)$
  \begin{align*}
    \mesh x < \delta \implies |R(f,x,t)- I| < \varepsilon.
  \end{align*}
  If $f$ is integrable on $[a,b]$, $I$ is the integral of $f$ over $[a,b]$ and it is denoted by
  \begin{align*}
    \int^b_a f(x) \,\mathrm{d} x = I = \lim_{\mesh x \to 0} \sum^n_{i=1} f(t_i) \Delta x_i.
  \end{align*}
\end{definition}

\begin{definition}[Upper Sum]
  Let $f:[a,b] \to \mathbb{R}$ be bounded above and let $x$ be a partition on $[a,b]$.
  The \emph{upper sum} of $f$ with respect to partition $x$ is
  \begin{align*}
    U(f, x) := \sum_{i=1}^n M_i \Delta x_i,
  \end{align*}
  where
  \begin{align*}
    M_i := \sup\left\{ f(x) \mid x \in [x_{i-1}, x_i]\right\}.
  \end{align*}
\end{definition}

\begin{definition}[Lower Sum]
  Let $f:[a,b] \to \mathbb{R}$ be bounded below and let $x$ be a partition on $[a,b]$.
  The \emph{lower sum} of $f$ with respect to partition $x$ is
  \begin{align*}
    L(f, x) := \sum_{i=1}^n m_i \Delta x_i,
  \end{align*}
  where
  \begin{align*}
    m_i := \inf\left\{ f(x) \mid x \in [x_{i-1}, x_i]\right\}.
  \end{align*}
\end{definition}

\begin{theorem}
  If $x^*$ is a refinement of $x$, then
  \begin{align*}
    L(f,x) &\leq L(f,x^*)\\
    U(f,x) &\geq U(f,x^*)
  \end{align*}
\end{theorem}
\begin{proof}
  Let $x^*$ be a refinement of partition $x$.
  Define
  \begin{align*}
    M_i &:= \sup \left\{f(x) \mid x \in [x_{i-1}, x_i]\right\},\\
    M_i^* &:= \sup \left\{f(x) \mid x \in [x^*_{i-1}, x^*_i]\right\},\\
    m_i &:= \inf \left\{f(x) \mid x \in [x_{i-1}, x_i]\right\},\\
    m_i^* &:= \inf \left\{f(x) \mid x \in [x^*_{i-1}, x^*_i]\right\}.
  \end{align*}
  Since $x^*$ is a refinement of partition $x$, there exists a sequence $p:[0,n]\cap\mathbb{N} \to [0,n^*]\cap\mathbb{N}$ such that $x = x^* \circ p$.
  Then $M$ and $m$ can be written as
  \begin{align*}
    M_i &= \sup \left\{ f(x) \mid x \in \left[x^*_{p_{i-1}}, x^*_{p_i}\right] \right\},\\
    m_i &= \inf \left\{ f(x) \mid x \in \left[x^*_{p_{i-1}}, x^*_{p_i}\right]\right\}.
  \end{align*}
  Since every subinterval $[x_{j-1}^*, x_j^*] \subseteq [x_{i-1}, x_i]$, we have
  \begin{align*}
    m_j^* \geq m_i \text{ and } M_j^* \leq M_i \text{ for all } j \in[p_{i-1}, p_i].
  \end{align*}
  So,
  \begin{align*}
    \sum_{j=p_{i-1}+1}^{p_i} m_j^* (x_j^* - x_{j-1}^*) &\geq \sum_{j=p_{i-1}+1}^{p_i}  m_i (x_j^* - x_{j-1}^*)\\
    \sum_{j=p_{i-1}+1}^{p_i} M_j^* (x_j^* - x_{j-1}^*) &\leq \sum_{j=p_{i-1}+1}^{p_i} M_i (x_j^* - x_{j-1}^*)
  \end{align*}
  Since $x^*$ is a refinement of $x$, the sum over $[x_{j-1}^*,x_{j_i}^*]$ for $j=p_{i-1}+1$ to $p_i$ telescopes to $[x_{i-1}, x_i]$.
  \begin{align*}
    \sum_{j=p_{i-1}+1}^{p_i} m_j^* (x_j^* - x_{j-1}^*) &\geq m_i (x_i - x_{i-1})\\
    \sum_{j=p_{i-1}+1}^{p_i} M_j^* (x_j^* - x_{j-1}^*) &\leq M_i (x_i - x_{i-1}).
  \end{align*}
  Summing over all intervals $[x_{i-1}, x_i]$
  \begin{align*}
    \sum_{i=1}^n \sum_{j=p_{i-1}+1}^{p_i} m_j^* (x_j^* - x_{j-1}^*) &\geq \sum_{i=1}^n  m_i (x_i - x_{i-1}),\\
    \sum_{i=1}^n \sum_{j=p_{i-1}+1}^{p_i} M_j^* (x_j^* - x_{j-1}^*) &\leq\sum_{i=1}^n   M_i (x_i - x_{i-1}).
  \end{align*}
  The double sum on the left spans all subintervals of $x^*$, so it equals the full lower or upper sum over $x^*$.
  The left side sum over sums is the full sum over the full interval $[x_0, x_{n^*}]$
  \begin{align*}
    \sum_{i=1}^{n^*=p_n} m_i^* (x_i^* - x_{i-1}^*) &\geq \sum_{i=1}^n  m_i (x_i - x_{i-1}),\\
    \sum_{i=1}^{n^*=p_n} M_i^* (x_i^* - x_{i-1}^*) &\leq\sum_{i=1}^n   M_i (x_i - x_{i-1}),
  \end{align*}
  which are the definitions of the upper and lower sums
  \begin{align*}
    L(f, x^*) &\geq L(f, x),\\
    U(f, x^*) &\leq U(f, x).
  \end{align*}
\end{proof}

\begin{theorem}
  Let $f:[a, b] \to \mathbb{R}$ be a bounded function. For any partition pair $(x, t)$
  \begin{align*}
    L(f, x) \leq R(f, x, t) \leq U(f, x).
  \end{align*}
\end{theorem}
\begin{proof}
  Define $M$ and $m$ such that
  \begin{align*}
    M_i &= \sup \{ f(x) \mid x \in [x_{i-1}, x_i] \},\\
    m_i &= \inf \{ f(x) \mid x \in [x_{i-1}, x_i] \}.
  \end{align*}
  Then for all pieces of the partition, the inequality holds
  \begin{align*}
    \forall i: m_i  \leq f(t_i) \leq M_i.
  \end{align*}
  So for the full partition, multiplying through by the size of the pieces
  \begin{align*}
    m_i \Delta x_i \leq f(t_i) \Delta x_i \leq M_i \Delta x_i.  
  \end{align*}
  Summing over all \( i \), we get:
  \begin{align*}
    \sum_{i=1}^n m_i \Delta x_i \leq \sum_{i=1}^n f(t_i) \Delta x_i \leq \sum_{i=1}^n M_i \Delta x_i.
  \end{align*}
  That is,
  \begin{align*}
    L(f, x) \leq R(f, x, t) \leq U(f, x).
  \end{align*}
\end{proof}

\begin{theorem}
  Let $f:[a,b] \to \mathbb{R}$ be a bounded function. 
  For all partitions, $x_1$, $x_2$,
  \begin{align*}
    L(f, x_1) \leq R(f, x^*, t) \leq U(f, x_2).
  \end{align*}
\end{theorem}
\begin{proof}
  Let $t^*$ be the common refinement of $t_1$ and $t_2$, then
  \begin{align*}
    L(f,t_1) \leq L(f,t^*)\\
    U(f,t^*) \leq U(f,t_2).
  \end{align*}
  By the above lemma
  \begin{align}
    L(f,t^*) \leq U(f,t^*)
  \end{align}
  Hence
  \begin{align}
    L(f,t_1) \leq U(f, t_2)
  \end{align}
\end{proof}

%\begin{corollary}
%The set $\{L(f,t): t\textrm{ is a partition} \}$ is bounded above and $\{U(f,t): t\textrm{ is a partition} \}$ is bounded below.
%\end{corollary}

\begin{theorem}
  Let $f: [a,b] \to \mathbb{R}$ be a bounded function. Then $f$ is integrable on $[a, b]$ if and only if, for every $\varepsilon>0$, there exists a partition $x$ of $[a, b]$ such that
  \begin{align*}
    U(f,x) - L(f,x) < \varepsilon.
  \end{align*}
\end{theorem}
\begin{proof}
  First, assume the bounded function $f$ is integrable. That is, there is an $I \in \mathbb{R}$ such that, for all $\varepsilon' > 0$ there is a $\delta > 0$ such that for every partition pair $(x, t)$
  \begin{align*}
    \mesh x < \delta \implies |R(f, x, t) - I| < \varepsilon'.
  \end{align*}
  Define $\varepsilon:=2\varepsilon'$, so that
  \begin{align*}
    |R(f, x, t) - I| < \frac{\varpsilon}{2}.
  \end{align*}
  If $f$ is integrable on $[a,b]$, $I$ is the integral of $f$ over $[a,b]$ and it is denoted by
  \begin{align*}
    \int^b_a f(x) \,\mathrm{d} x = I = \lim_{\mesh x \to 0} \sum^n_{i=1} f(t_i) \Delta x_i.
  \end{align*}
\end{proof}


\section{Primitives}

\begin{definition}[Primitive]
  Let $f$ be continuous on $(a,b)$. The function $F$ is a primitive of $f$, if $F$ is differentiable on $(a,b)$ and
  \begin{align}
    \forall x\in(a,b): F'(x) = f(x)
  \end{align}

\end{definition}

\begin{lemma}
  If $f$ is continuous on $[a,b]$, then the function $F(x) = \int^x_a f(t)\, \dif t$ is uniformly continuous on $[a,b]$.
\end{lemma}

\begin{lemma}
  If $f$ is continuous on $[a,b]$, then the function $F(x) = \int^x_a f(t)\, \dif t$ is differentiable on $(a,b)$.
\end{lemma}

\begin{theorem}
  If $f$ is continuous on $[a,b]$, then the function $F(x) = \int^x_a f(t)\, \dif t$ is a primitive of $f$.
\end{theorem}


\begin{theorem}
  If $f$ is integrable on $[a,b]$, $F$ is differentiable on $[a,b]$, and $F$ is a primitive of $f$, then
  \begin{align}
    \int^b_a f(x) \, \dif x = F(b) - F(a)
  \end{align}
\end{theorem}



\partmark{Functional Analysis}
\newpage
\section{Normed Spaces}
\section{Functionals}
\section{Fr\'{e}chet \& Gateaux Derivatives}
\section{Weak Convergence}
\end{document}
